\documentclass[11pt,a4paper]{article}
\usepackage[margin=24mm]{geometry}
\usepackage[T1]{fontenc}
\usepackage{lmodern,amsmath,amssymb,amsthm,booktabs,microtype,dsfont,braket,xurl,longtable}
\usepackage[colorlinks=true,linkcolor=blue,urlcolor=blue]{hyperref}
\usepackage{eso-pic,embedfile,tikz-cd}
\newcommand{\A}{\mathcal A}
\newcommand{\Tr}{\operatorname{Tr}}
\newcommand{\id}{\mathds{1}}
\newcommand{\D}{\mathcal D}
\newtheorem{theorem}{Theorem}
\newtheorem{lemma}{Lemma}
\newcommand{\HS}{\mathrm{HS}}
\newcommand{\diag}{\operatorname{diag}}
\newcommand{\E}{\mathbb E}
\DeclareMathOperator{\Sym}{Sym}
\setlength{\emergencystretch}{2em}
\setcounter{tocdepth}{2}
\title{Monotonicity of total anticoherence\\under $\mathrm{SU}(2)$-covariant channels\\[3mm]\large Fidelity $\mathcal A_t^{T,F}$ and von Neumann $\mathcal A_t^{T,\mathrm{vN}}$ anticoherence measures at orders $t\le4$\\Uniform proofs with embedded exact data}
\author{}
\date{23 November 2026}
\hypersetup{pdftitle={Monotonicity of fidelity and von Neumann total anticoherence under SU(2)-covariant channels at orders t <= 4}}
\begin{document}
\embedfile[filespec=exact-data.zip,desc={Exact rational data and checking programs for orders two, three and four},mimetype=application/zip]{exact-data.zip}
\maketitle
\noindent We prove fidelity monotonicity for every input density matrix at
orders $t=1,2,3,4$ by comparing shifted determinants. Von Neumann
monotonicity then follows from a single implication valid at every order.
Section~\ref{sec:lean} identifies the Lean theorems for both measures.

\medskip
\noindent The exact data and verification scripts used in the
computer-assisted proofs for $t=3,4$ are provided in \texttt{exact-data.zip},
embedded in this PDF and available in the repository.
See Section~\ref{sec:data} for verification instructions.
\clearpage
{\small\tableofcontents}
\clearpage
\section{Statement in the manuscript's notation}
Put $N=2j$, let $\mathcal H_j$ carry the irreducible spin-$j$
representation, and identify it with $\operatorname{Sym}^N(\mathbb C^2)$.
For $\rho\in\D(\mathcal H_j)$, write $\rho_S$ for the corresponding
symmetric $N$-qubit state and
\[
 \rho_t=\Tr_{N-t}(\rho_S),\qquad
 \rho_0=\frac{\id_{N+1}}{N+1},\qquad
 \rho_0^{(t)}=\frac{\id_{t+1}}{t+1}.
\]
Every reduced trace and determinant is taken on the symmetric subspace
$\mathcal H_{t/2}$, of dimension $d=t+1$. Let $\rho'=\Phi(\rho)$
and $\rho'_t=\Tr_{N-t}(\rho'_S)$; the prime does not denote an
independently defined channel on all $t$-qubit states.

The manuscript uses the squared Uhlmann--Jozsa fidelity convention,
$F(\rho_t,\rho_0^{(t)})=\frac{(\Tr\sqrt{\rho_t})^2}{t+1}$.
The two anticoherence measures are
\begin{align}
 \A_t^{T,F}(\rho)&=\frac{(\Tr\sqrt{\rho_t})^2-1}{t},\label{eq:F}\\
 \A_t^{T,\mathrm{vN}}(\rho)&=
 -\frac{\Tr(\rho_t\log\rho_t)}{\log(t+1)}.\label{eq:VN}
\end{align}
Write $S(X)=-\Tr(X\log X)$, with $0\log0=0$ and natural logarithms.

\begin{theorem}\label{thm:main}
Let $t\in\{1,2,3,4\}$ and $N\ge t$. For every completely positive,
trace-preserving channel $\Phi$ on $\mathcal B(\mathcal H_{N/2})$
satisfying
\[
 \Phi(U\rho U^\dagger)=U\Phi(\rho)U^\dagger
\]
for every spin-$N/2$ rotation $U$ and every density matrix $\rho$,
one has
\[
 \boxed{\A_t^{T,F}(\Phi(\rho))\ge\A_t^{T,F}(\rho),\qquad
 \A_t^{T,\mathrm{vN}}(\Phi(\rho))\ge\A_t^{T,\mathrm{vN}}(\rho).}
\]
The statement includes all density matrices, including pure states.
\end{theorem}

\paragraph{Uniform proof structure.}
Fix $N$, $t$ and $\Phi$ as in Theorem~\ref{thm:main}. Every comparison in
the diagram is for \emph{all} input density matrices of this same channel.
In the shifted determinant comparison, $s$ ranges over all positive reals.
\[
\begin{tikzcd}[column sep=2em,row sep=2.1em]
 \det\rho'_t\ge\det\rho_t \arrow[r,Rightarrow,"{\S\ref{sec:shift}}"] &
 \det(s\id_d+\rho'_t)\ge\det(s\id_d+\rho_t)
    \arrow[d,Rightarrow,"{\S\ref{sec:spectral}}"] & \\
 & \Tr\sqrt{\rho'_t}\ge\Tr\sqrt{\rho_t}
    \arrow[r,Rightarrow,"{\S\ref{sec:fidelity-implies-entropy}}"] &
 S(\rho'_t)\ge S(\rho_t)
\end{tikzcd}
\]
The last arrow uses fidelity monotonicity on the entire segment from
each $\rho$ to the maximally mixed spin state $\rho_0$. Only the first
inequality requires a proof specific to the order.
\begin{center}
\begin{tabular}{@{}cccl@{}}
\toprule
$t$&$d=t+1$&Proper-order domain&Determinant certificate\\
\midrule
1&2&$N\ge2$&Bloch-vector determinant identity\\
2&3&$N\ge3$&Physical-marginal polynomial argument\\
3&4&$N\ge4$&Four base channels and the all-spin hull\\
4&5&$N\ge5$&Fourteen base channels and the all-spin hull\\
\bottomrule
\end{tabular}
\end{center}


Sections~\ref{sec:t1}--\ref{sec:t4} establish
\begin{equation}\label{eq:det}
 \det\rho'_t\ge\det\rho_t
\end{equation}
for $N>t$. At $N=t$ it follows from unitality and trace preservation
(Section~\ref{sec:prelim}). The exact data for $t=3,4$ are embedded in
this PDF (Section~\ref{sec:data}). For $t=5$ and $N>5$, see the
\href{https://github.com/aenictusGitHub/mixed-spin-anticoherence/tree/main/su2-order-five-monotonicity}{order-five supplement}.

\section{Common algebraic and spectral preliminaries}\label{sec:prelim}
This section supplies the channel and matrix arguments used at all four
orders. Ordinary finite-dimensional linear algebra, the spectral theorem,
and elementary single-variable calculus are the background conventions.
All finite certificate objects used later are defined by the exact data
embedded in this PDF; their interpretation and checking rules are given
in Section~\ref{sec:data}. They are part of this document, not references
to a separate repository or a Lean computation.

\subsection{Symmetric reduction and its transitivity}
Let $|n,a\rangle$, $0\le a\le n$, be the normalized symmetric sum of all
binary words of length $n$ having $a$ entries equal to one. Splitting a
word after its first $t$ entries gives the isometric embedding
\[
 |N,a\rangle\longmapsto
 \sum_{r+b=a}\sqrt{\frac{\binom tr\binom{N-t}b}{\binom Na}}
 |t,r\rangle\otimes|N-t,b\rangle.
\]
Its coefficients have squared sum one by the binomial identity
$\sum_r\binom tr\binom{N-t}{a-r}=\binom Na$; different $a$ have
orthogonal images. Thus, with out-of-range entries interpreted as zero,
\begin{equation}\label{eq:reduction-kraus}
 R_{N,t}(X)=\sum_{b=0}^{N-t}F_bXF_b^\dagger,
 \qquad (F_b)_{r,a}=\mathbf1_{a=r+b}
 \sqrt{\frac{\binom tr\binom{N-t}b}{\binom Na}}.
\end{equation}
The binomial identity also gives $\sum_bF_b^\dagger F_b=I$.
This proves complete positivity and trace preservation directly.
Taking partial traces in two stages proves
$R_{m,t}R_{N,m}=R_{N,t}$ for $t\le m\le N$.
The embedding is the restriction of the ordinary tensor product of
qubit rotations, so reduction intertwines rotations. An invariant
operator on spin $t/2$ is scalar: commuting with $J_z$ makes it diagonal,
and commuting with $J_+$ makes all diagonal entries equal.
Consequently $R_{N,t}(I_{N+1}/(N+1))=I_{t+1}/(t+1)$.

\subsection{Tensor ranks and the covariant-channel simplex}
In the magnetic basis $|j,m\rangle$,
\[
 J_z|j,m\rangle=m|j,m\rangle,\qquad
 J_+|j,m\rangle=\sqrt{(j-m)(j+m+1)}|j,m+1\rangle,
 \quad J_-=J_+^\dagger.
\]
The matrix space is the orthogonal direct sum of ranks $L=0,\ldots,N$.
One direct construction starts each rank at $J_+^L$ and applies
$[J_-,\cdot]$ successively. Its $2L+1$ weights are $L,L-1,\ldots,-L$;
their dimensions sum to the matrix-space dimension in the form
$\sum_{L=0}^N(2L+1)=(N+1)^2$. The Casimir of the commutator action
has eigenvalue $L(L+1)$ on rank $L$, proving orthogonality between ranks.
Normalize these vectors in Hilbert--Schmidt norm to obtain the manuscript's
$T_{LM}$, up to immaterial phase choices. Equivalently,
\[
 (T_{LM})_{m,n}=\sqrt{\frac{2L+1}{2j+1}}
 \langle j,n;L,M\mid j,m\rangle.
\]
Each rank has a unique highest-weight line: an operator commuting with
$J_+$ is determined by its value on the lowest-weight vector, and the
weight equation fixes that value. A covariant linear map must preserve
these highest-weight lines and commute with lowering, and therefore
acts by a scalar $f_L$ on rank $L$. Hermiticity makes $f_L$ real;
trace preservation fixes $f_0=1$. Reduction intertwines the same ranks,
so its nonzero rank-$L$ components acquire exactly $f_L$, for $L\le t$.
The reduction coefficients on these ranks are nonzero: applying the
adjoint of reduction to $J_+^L$ gives a positive nonzero multiple of
the global $J_+^L$, by expanding the collective raising operator into
products on distinct qubits. Thus matching the multipliers matches
the entire reduced linear map.

Here is a direct proof that the admissible full channels form a simplex.
Write $d_N=N+1$, $|\Omega\rangle=\sum_a|a,a\rangle$ and
$C_\Phi=(\Phi\otimes\mathrm{id})(|\Omega\rangle\langle\Omega|)$.
Complete positivity implies $C_\Phi\succeq0$, and trace preservation
is equivalent to $\operatorname{Tr}_{\rm out}C_\Phi=I$.
Conversely a positive $C=\sum_s|v_s\rangle\langle v_s|$ gives Kraus
matrices by reshaping $v_s$, and hence gives a completely positive map;
the partial-trace condition is exactly the Kraus normalization.
Covariance says that $C_\Phi$ commutes with $U\otimes\bar U$.
Time reversal identifies the conjugate spin with spin $j$, so this
representation decomposes into one copy of each spin $K=0,\ldots,N$.
For completeness, the tensor-product weight multiplicity at weight
$K\ge0$ is $N+1-K$; subtracting its value at $K+1$ gives one
highest-weight vector of spin $K$. Applying lowering generates that
copy, and the dimensions sum to $d_N^2$.
Let $P_K$ denote its orthogonal projector. Since the copies are
inequivalent, a commuting Hermitian operator is scalar on each copy:
commutation with weights and raising proves this within a copy and
excludes maps between distinct highest weights. Thus
\[
 C_\Phi=\sum_{K=0}^N c_KP_K,\qquad c_K\ge0.
\]
The partial trace of $P_K$ is invariant, and its trace is $2K+1$;
therefore it is $(2K+1)I/d_N$. Put
$p_K=c_K(2K+1)/d_N$. Trace preservation gives $\sum_Kp_K=1$, and
\begin{equation}\label{eq:channel-simplex}
 \Phi=\sum_{K=0}^Np_K\Phi_K^{(N)},\qquad
 C_{\Phi_K^{(N)}}=\frac{d_N}{2K+1}P_K.
\end{equation}
This derives the simplex from the physical assumptions.

\subsection{An explicit computation of the four relevant multipliers}
Put $D=N(N+2)$, $u=K(K+1)$, and define
\[
 B_r(u)=\prod_{i=0}^{r-1}(u-i(i+1)),\quad
 D_r=\prod_{i=0}^{r-1}(D-i(i+2)),\quad B_0=D_0=1.
\]
For $0\le L\le4$ and $N\ge L$, the multipliers are
\begin{equation}\label{eq:racah-derived}
 f_L(N,K)=\sum_{r=0}^L(-1)^r\binom Lr\binom{L+r}r\frac{B_r(u)}{D_r}.
\end{equation}
All denominators are positive when $N\ge L$. We give an algebraic way to
obtain and check this expression, avoiding an appeal to a table of
recoupling coefficients.

Set $X=-\sum_aJ_a\otimes J_a^{\mathsf T}$ and
$Z_L=\sum_{M=-L}^LT_{LM}\otimes\overline{T_{LM}}$.
The total-spin Casimir of $U\otimes\bar U$ is $2j(j+1)I+2X$;
on $P_K$, $X$ therefore has eigenvalue $u/2-j(j+1)$.
We compute $Z_L$ directly, including the finite identities that identify
its eigenvalues with the physical channel multipliers.

Use a real diagonal change of basis on spin $N/2$ so that, in ascending
weight order $k=0,\ldots,N$, the raising and lowering matrices become
\[
 E|k\rangle=(N-k)|k+1\rangle,\qquad
 F|k\rangle=k|k-1\rangle,\qquad
 J_z|k\rangle=(k-N/2)|k\rangle.
\]
This change is a similarity; Hilbert--Schmidt norms below are always
computed before the change, in the orthonormal physical basis.
Let $W_{L,r}=(\operatorname{ad}F)^rE^L$. It has just one nonzero band:
\[
 W_{L,r}|k\rangle=w_{L,r}(k)|k+L-r\rangle,
\]
where ordinary matrix multiplication gives the complete recurrence
\begin{align*}
 w_{L,0}(k)&=\prod_{i=0}^{L-1}(N-k-i),\\
 w_{L,r+1}(k)&=(k+L-r)w_{L,r}(k)-k w_{L,r}(k-1).
\end{align*}
Its physical norm at $r=0$ is
\[
 s_L=\operatorname{Tr}(J_-^LJ_+^L)
 =(L!)^2\sum_{k=0}^{N-L}\binom{k+L}{L}\binom{N-k}{L}
 =(L!)^2\binom{N+L+1}{2L+1}.
\]
For the last equality, count $(2L+1)$-element subsets of
$\{0,\ldots,N+L\}$ according to their middle element $k+L$.
Since $(\operatorname{ad}J_-)^*=\operatorname{ad}J_+$ in the
Hilbert--Schmidt inner product, the spin-$L$ ladder relations give
\[
 \|W_{L,r}\|_{\rm HS}^2
 =s_L\,r!\frac{(2L)!}{(2L-r)!}.
\]
Indeed, raising the $r$th descendant gives
$r(2L-r+1)$ times the $(r-1)$st descendant, which proves the formula
inductively. Dividing by these norms gives the real tensors $T_{L,L-r}$
(up to signs that cancel in $Z_L$).

Apply the same real similarity on both tensor factors. The coefficient
of $|k+a,h+a\rangle$ in $Z_L|k,h\rangle$, for $a=L-r$, is therefore
\[
 z_{L,a}(k,h)=\frac{w_{L,r}(k)w_{L,r}(h)}
 {s_L\,r!(2L)!/(2L-r)!}.
\]
There are no other bands. In these coordinates $X$ has three bands,
with coefficients
\[
 x_0(k,h)=-(k-N/2)(h-N/2),\quad
 x_1(k,h)=-\frac{(N-k)(N-h)}2,\quad
 x_{-1}(k,h)=-\frac{kh}2.
\]
To compute every power needed below, set $x_a^{(0)}=\mathbf1_{a=0}$
and use
\[
 x_a^{(p+1)}(k,h)=\sum_{b=-1}^1x_{a-b}^{(p)}(k,h)
 x_b(k+a-b,h+a-b).
\]
All coefficients outside the finite basis vanish at the boundary by
the raising and lowering formulas.

Put $c=N(N+2)/4$. Substitution of the two band recurrences gives
$d_NZ_L/(2L+1)=P_L(X)$, where $P_0=1$ and
\begin{align*}
 P_1(x)&=-x/c,\\
 P_2(x)&=\frac{6x^2+3x-2c^2}{c(4c-3)},\\
 P_3(x)&=\frac{-10x^3-20x^2+(6c^2-4c-6)x+5c^2}
 {c(c-2)(4c-3)},\\
 P_4(x)&=\frac{70x^4+350x^3+(-60c^2+100c+390)x^2}
 {c(c-2)(4c-15)(4c-3)}\\
 &\quad+\frac{(-170c^2+120c+90)x+6c^4-24c^3-81c^2}
 {c(c-2)(4c-15)(4c-3)}.
\end{align*}
Here is a precise coefficient test for these four identities: for each
$L=1,2,3,4$ and each $a=-L,\ldots,L$, expand $P_L(x)=\sum_pp_{L,p}x^p$
and subtract
\[
 \sum_{p=0}^Lp_{L,p}x_a^{(p)}(k,h)
 -\frac{N+1}{2L+1}z_{L,a}(k,h).
\]
Each numerator is identically zero as a polynomial in $N,k,h$ after
clearing the displayed denominators. Thus there are respectively
$3,5,7,9$ rational identities, specified completely by the recurrences
above. This is a symbolic identity check for all $N\ge L$, not a test
at selected spins. The denominators are nonzero on this domain.

Finally the Choi trace pairing gives
\[
 \operatorname{Tr}(C_\Phi Z_L^\dagger)
 =\sum_M\operatorname{Tr}\bigl(T_{LM}^\dagger\Phi(T_{LM})\bigr)
 =(2L+1)f_L.
\]
For the vertex $C_\Phi=d_NP_K/(2K+1)$, the polynomial identity just
proved makes the left side $(2L+1)P_L(u/2-c)$. Hence
$f_L(N,K)=P_L(u/2-c)$. Substitution into the four displayed polynomials
and collecting powers of $u$ gives exactly~\eqref{eq:racah-derived}.
This derives the multiplier formulas from the Choi definition and
finite matrix operations, without assuming a recoupling formula.

\subsection{Rectangular channels and how their coefficients are defined}
All positive auxiliary maps used below have a concrete construction.
For spins $a,b,j$ with $|a-b|\le j\le a+b$ and integral $a+b-j$,
construct the normalized spin-$j$ highest-weight vector in
$\mathcal H_a\otimes\mathcal H_b$ by solving
\[
 (J_z\otimes I+I\otimes J_z)v=jv,\qquad
 (J_+\otimes I+I\otimes J_+)v=0,\qquad\|v\|=1.
\]
Choose its first nonzero coefficient positive. The raising equation is
a first-order recurrence in the coefficients, so determines them up
to scale. Obtain the other columns by lowering:
\[
 V|j,m-1\rangle=
 \frac{(J_-\otimes I+I\otimes J_-)V|j,m\rangle}
 {\sqrt{(j+m)(j-m+1)}}.
\]
The commutation relations prove by induction that these columns are
orthonormal and intertwine all three spin generators. Exponentiating
proves rotation covariance. Their entries are the real
Clebsch--Gordan coefficients used throughout this document.
The channel $G(\rho)=\operatorname{Tr}_b(V\rho V^\dagger)$ therefore
has Kraus entries
\[
 (A_\mu)_{m,n}=\langle a,m;b,\mu\mid j,n\rangle,
 \qquad G(\rho)=\sum_{\mu=-b}^b A_\mu\rho A_\mu^\dagger.
\]
It is completely positive and trace preserving, and sends the invariant
input state to $I_{2a+1}/(2a+1)$. The choice $j=N/2$, $a=t/2$,
$b=(N-t)/2$ recovers~\eqref{eq:reduction-kraus}.
The choice $a=j,b=K$ gives the vertex in~\eqref{eq:channel-simplex}:
the vectorized Kraus span is its unique rank-$K$ copy and normalization
fixes its Choi scalar. Equivalent adjoint/phase choices in the square
Kraus convention give the same vertex.

Every finite auxiliary multiplier table in this document can therefore
be evaluated without an external convention. Compute $G(T_{L0}^{(j)})$
by the displayed Kraus sum and divide its $T_{L0}^{(a)}$ coefficient
by the corresponding coefficient of $R_{N,t}(T_{L0}^{(j)})$.
Covariance makes this ratio valid for every $M$. The embedded tables
and checkers use exactly these finite sums.

\subsection{Spin reversal, positive maps, and determinant mixtures}
In magnetic order define $V|j,m\rangle=(-1)^{j-m}|j,-m\rangle$ and
$\Theta(X)=VX^{\mathsf T}V^\dagger$. Transpose preserves positivity
and spectrum, so $\Theta$ does too; it is unital and trace preserving.
The generator formulas show covariance and
$\Theta(T_{LM})=(-1)^LT_{LM}$ on Hermitian rank components.
For half-integral spin, $V^{\mathsf T}=-V$, whence
$\langle\psi,V\bar\psi\rangle=0$. This proves the positivity of
$X\mapsto[\operatorname{Tr}(X)I-X-\Theta(X)]/(d-2)$ for even $d>2$
by first checking rank-one projectors and then summing their spectral
decompositions. Its trace and identity image give trace preservation
and unitality.

If $\Psi$ is any positive unital trace-preserving map on $d\times d$
matrices, let $p_i$ be the input eigenvalues, $P_i$ their rank-one
spectral projectors, and $v_j$ an orthonormal output eigenbasis.
Then the output eigenvalues satisfy
\[
 q_j=\sum_iD_{ji}p_i,\qquad
 D_{ji}=\langle v_j,\Psi(P_i)v_j\rangle\ge0,
 \quad \sum_iD_{ji}=\sum_jD_{ji}=1.
\]
For positive definite input, scalar concavity of $\log$ gives
$\sum_j\log q_j\ge\sum_i\log p_i$. This proves determinant increase.
For singular input it follows simply from output positivity.
This also proves the assertion for $N=t$, since the full channel is
unital by covariance and irreducibility.

Finally log determinant is concave on positive definite matrices.
Indeed, along $A(s)=A+sH\succ0$ its first and second derivatives are
\[
 \frac{d}{ds}\log\det A(s)=\operatorname{Tr}(A(s)^{-1}H),\qquad
 \frac{d^2}{ds^2}\log\det A(s)
 =-\operatorname{Tr}\bigl[(A(s)^{-1/2}HA(s)^{-1/2})^2\bigr]\le0.
\]
The first identity follows by expanding $\det(I+hA^{-1}H)$ to first
order; differentiating $A^{-1}A=I$ gives the second.
Thus, if positive outputs $\tau_i$ each have determinant at least that
of a positive definite input $\sigma$, then
\[
 \log\det\Bigl(\sum_iw_i\tau_i\Bigr)
 \ge\sum_iw_i\log\det\tau_i\ge\log\det\sigma
 \quad(w_i\ge0,\ \sum_iw_i=1).
\]
For singular $\sigma$, positivity suffices again. This is the precise
convex-mixture argument used both for channel vertices and for the
available-extension hulls.

\clearpage
\section{Determinant comparison at order one}\label{sec:t1}
Let $N\ge1$, $\sigma=\rho_1$ and $\tau=\rho'_1$. We prove
$\det\tau\ge\det\sigma$ for every physical input.

\subsection{The rank-one multiplier and the physical reduction}
Permutation symmetry and
$\mathbf J=\tfrac12\sum_{a=1}^N\boldsymbol\sigma^{(a)}$ give
\[
 \sigma=\frac{I_2+\mathbf r\cdot\boldsymbol\sigma}{2},\qquad
 \mathbf r=\frac{\langle\mathbf J\rangle_\rho}{j},\qquad
 r=|\mathbf r|\le1,
\]
where $j=N/2$ and $\boldsymbol\sigma$ is the vector of Pauli matrices.
By the rank decomposition in Section~\ref{sec:prelim}, covariance
multiplies every rank-one component by the same real number $f_1$.
Reduction intertwines rotations, so
\[
 \tau=\frac{I_2+f_1\mathbf r\cdot\boldsymbol\sigma}{2}.
\]
To bound $f_1$, apply this formula to a coherent full input, whose
one-qubit reduction has $r=1$. The output eigenvalues
$(1\pm f_1)/2$ must be nonnegative. Hence $|f_1|\le1$ for the channel,
independently of the chosen input.

\subsection{Determinant monotonicity}
The Pauli identity $(\mathbf r\cdot\boldsymbol\sigma)^2=r^2I_2$
gives the eigenvalues $(1\pm r)/2$ of $\sigma$. Therefore
\[
 \det\sigma=\frac{1-r^2}{4},\qquad
 \det\tau=\frac{1-f_1^2r^2}{4},\qquad
 \boxed{\det\tau-\det\sigma=\frac{(1-f_1^2)r^2}{4}\ge0.}
\]
No invertibility assumption was used: this includes pure reductions
($r=1$), the maximally mixed reduction ($r=0$), and arbitrary mixed
full inputs.

\clearpage
\section{Determinant comparison at order two}\label{sec:t2}
Let $N\ge3$, $\sigma=\rho_2$ and $\tau=\rho'_2$. We prove
$\det\tau\ge\det\sigma$ for every physical input. The following argument
includes the channel-multiplier bounds, the physical extension constraint,
and both scalar estimates used in the determinant identity.

\subsection{The two channel multipliers}
Use the Hilbert--Schmidt orthonormal tensor operators $T_{LM}$ of the
manuscript and put
\[
 f=f_1,\qquad g=f_2,\qquad S=N(N+2),\qquad c=j(j+1)=S/4.
\]
We give an elementary justification that these two real multipliers exist;
no decomposition of the Choi representation is needed.
Let $H=J_z$, $E=J_x+iJ_y$, $E^\dagger=J_x-iJ_y$. In the usual
weight basis $\ket{k}$, $0\le k\le N$,
\[
 H\ket{k}=(k-j)\ket{k},\qquad
 E\ket{k}=\sqrt{(k+1)(N-k)}\ket{k+1}.
\]
State covariance extends to all operators by complex linearity: positive
operators are nonnegative multiples of states, and every matrix is a
complex linear combination of positive operators. Differentiating
rotation covariance shows that $\Phi$ commutes with
$[J_a,\cdot]$, hence with $[H,\cdot]$, $[E,\cdot]$, and
$[E^\dagger,\cdot]$.
If $[H,X]=LX$ and $[E,X]=0$, then $X\ket{0}$ is a multiple of
$\ket{L}$. Moreover $X E^k\ket{0}=E^kX\ket{0}$ determines $X$
on the entire basis. Thus $X$ is a multiple of $E^L$ for
$L=0,1,2$. Applying this to $X=\Phi(E^L)$ shows that $\Phi$
is scalar on each space generated from $E^L$ by repeated commutators
with $E^\dagger$. For $L=0$ this is the scalar space; for $L=1$
it is the span of $J_a$; for $L=2$ it is the span of
\[
 Q_{ab}=\frac{J_aJ_b+J_bJ_a}{2}-\frac c3\delta_{ab}\id.
\]
For example,
\[
 [E^\dagger,E]=-2H,\quad
 [E^\dagger,E^2]=-2(HE+EH),\quad
 [E^\dagger,[E^\dagger,E^2]]=12H^2-4c\id;
\]
these identities, their adjoints, and $\sum_aQ_{aa}=0$ give all the
listed components.
Hermiticity preservation makes the scalars real, and trace preservation
makes the rank-zero scalar one. Hence
\[
 \Phi(\id)=\id,\qquad \Phi(J_a)=fJ_a,\qquad \Phi(Q_{ab})=gQ_{ab}.
\]
The trace adjoint also commutes with these commutators and preserves
Hermiticity. The same argument applies to it, and trace pairing with
any nonzero Hermitian component shows that its rank-one and rank-two
scalars are the same $f,g$. This also justifies the transformation of
expectation values used below.

\begin{lemma}[Multiplier bounds]\label{t2:lem:multipliers}
Every channel in Theorem~\ref{thm:main} satisfies
\begin{align}
 -\frac23\le g\le1,\qquad 3|f|\le2+g,\label{t2:eq:region}\\
 3Sf^2-6f-S\le2(S-3)g.\label{t2:eq:curved}
\end{align}
Consequently $|f|\le1$ and $|g|\le1$; if $g<-1/2$, then $|f|<1/3$.
\end{lemma}

\begin{proof}
Set $d=N+1$ and let $C_\Phi\succeq0$ be the unnormalized Choi
matrix, with convention
\begin{equation}\label{t2:eq:choi-trace}
 \Tr[C_\Phi(B\otimes C^{\mathsf T})]=\Tr[B\Phi(C)].
\end{equation}
Complete positivity gives $C_\Phi\succeq0$ by applying
$\Phi\otimes\mathrm{id}$ to the unnormalized maximally entangled
projector. Trace preservation gives $\Tr C_\Phi=d$.
Thus $\omega=C_\Phi/d$ is a density matrix.
Define the Hermitian observable
\[
 X=-\sum_{a=1}^3J_a\otimes J_a^{\mathsf T}.
\]
Only the following two polynomial identities are required:
\begin{equation}\label{t2:eq:quadrupole}
 \sum_{a,b}Q_{ab}\otimes Q_{ab}^{\mathsf T}
 =X^2+\frac X2-\frac{c^2}3\id,
 \qquad
 \sum_{a,b}Q_{ab}^2=\frac{c(4c-3)}6\id.
\end{equation}
For the first identity, put $B_{ab}=(J_aJ_b+J_bJ_a)/2$ and expand:
\[
 \sum_{a,b}B_{ab}\otimes B_{ab}^{\mathsf T}
 =X^2+\frac12\sum_{a,b}J_aJ_b\otimes[J_a,J_b]^{\mathsf T}
 =X^2+\frac X2.
\]
The last equality uses $\sum_{a,b}\epsilon_{abk}J_aJ_b=iJ_k$;
subtracting the scalar trace parts gives the first identity.
For the second, expansion of
$\sum_b[J_b,[J_b,J_a]]=2J_a$ gives
$\sum_bJ_bJ_aJ_b=(c-1)J_a$. Thus, for
$B_{ab}=(J_aJ_b+J_bJ_a)/2$,
\[
 \sum_{a,b}B_{ab}^2
 =\frac12\sum_{a,b}(J_aJ_bJ_aJ_b+J_aJ_bJ_bJ_a)
 =\left(c^2-\frac c2\right)\id,
\]
which proves the second identity.
The trace pairing and the two multipliers now give
\[
 \Tr(\omega X)=-cf,\qquad
 \Tr(\omega X^2)=\frac{c(4c-3)}6g+\frac{cf}{2}+\frac{c^2}{3}.
\]
Nonnegative variance,
$\Tr(\omega X^2)\ge[\Tr(\omega X)]^2$, is precisely
\eqref{t2:eq:curved} after multiplication by $12/c$.

The affine marginal law \eqref{t2:eq:transform}, established in the next
section, applied to a spin-coherent input gives reduced output
eigenvalues
\[
 \frac{2+g+3f}{6},\qquad \frac{2+g-3f}{6},\qquad\frac{1-g}{3}.
\]
Their nonnegativity gives $g\le1$ and $3|f|\le2+g$.
Finally, if $g<-2/3$, \eqref{t2:eq:curved} and $S\ge15$ would give
\[
 0>9(Sf-1)^2+S^2-12S-9
   =9(Sf-1)^2+(S-15)(S+3)+36\ge0.
\]
This proves \eqref{t2:eq:region}. If $g<-1/2$,
\eqref{t2:eq:curved} implies $Sf^2-2f-1<0$.
Since $S\ge15$, this forces $-1/3<f<1/3$: otherwise
$Sf^2-2f-1\ge15f^2-2f-1=(3f-1)(5f+1)\ge0$.
\end{proof}

\subsection{Physical two-qubit reductions}
Write the symmetric two-qubit reduction in its Pauli form
\begin{equation}\label{t2:eq:pauli}
 \rho_2=\frac14\left[\id\otimes\id+
 \sum_{a=1}^3r_a(\sigma_a\otimes\id+\id\otimes\sigma_a)
 +\sum_{a,b=1}^3T_{ab}\sigma_a\otimes\sigma_b\right].
\end{equation}
Here $\mathbf r\in\mathbb R^3$, $T=T^{\mathsf T}\in M_3(\mathbb R)$,
and $\Tr T=1$, since $\sum_a\sigma_a\otimes\sigma_a=2\mathsf F-\id$
and the swap $\mathsf F$ is the identity on the symmetric subspace.
In the orthonormal Cartesian spin-one basis
\[
 \ket{x}=\frac{\ket{11}-\ket{00}}{\sqrt2},\quad
 \ket{y}=\frac{i(\ket{11}+\ket{00})}{\sqrt2},\quad
 \ket{z}=\frac{\ket{01}+\ket{10}}{\sqrt2},
\]
put $A=\id_3-T$ and let $[\mathbf r]_\times\mathbf v=
\mathbf r\times\mathbf v$. Substitution in \eqref{t2:eq:pauli} gives
\begin{equation}\label{t2:eq:cartesian}
 2\rho_2=A+i[\mathbf r]_\times,\qquad
 8\det\rho_2=\det A-\mathbf r^{\mathsf T}A\mathbf r.
\end{equation}
To check the determinant identity, diagonalize $A$ by a proper real
rotation. The matrix then has diagonal $a,b,c$ and off-diagonal entries
$-ir_3,ir_2,-ir_1$ and their conjugates; its determinant is
$abc-ar_1^2-br_2^2-cr_3^2$.

There is a useful bound specific to reductions of symmetric $N$-qubit
states. With $J_a=\frac12\sum_{n=1}^N\sigma_a^{(n)}$, permutation
symmetry gives
\[
 \langle J_a\rangle=\frac N2r_a,\qquad
 \left\langle\frac{J_aJ_b+J_bJ_a}{2}\right\rangle
 =\frac N4\delta_{ab}+\frac{N(N-1)}4T_{ab}.
\]
The state has a symmetric three-qubit reduction with the same two-qubit
marginal. Applying these identities to that three-qubit state and using
nonnegative variance of its collective spin gives
$\id+2T-3\mathbf r\mathbf r^{\mathsf T}\succeq0$.
Hence $\id+2T\succeq0$. Testing \eqref{t2:eq:cartesian} on real
vectors also gives $A\succeq0$. Therefore
\begin{equation}\label{t2:eq:A-bounds}
 0\preceq A\preceq\frac32\id_3,\qquad\Tr A=2.
\end{equation}
This uses only the symmetric three-qubit extension, available for every
$N\ge3$.

Write $\rho'=\Phi(\rho)$. The partial trace intertwines rotations,
so the rank-one and rank-two parts of the reduction acquire the same
multipliers $f_1=f$ and $f_2=g$ as the global tensor components.
Equivalently, the above first and second moments give directly
\begin{equation}\label{t2:eq:transform}
 \mathbf r'=f\mathbf r,\quad T'=gT+\frac{1-g}{3}\id_3,
 \quad A'=\alpha\id_3+gA,\qquad\alpha=\frac{2(1-g)}3.
\end{equation}

\subsection{Determinant monotonicity}
We prove two elementary estimates for $A$ and $A'$.

\begin{lemma}\label{t2:lem:volume}
Under \eqref{t2:eq:A-bounds} and \eqref{t2:eq:region}, $\det A'\ge\det A$.
\end{lemma}
\begin{proof}
Write the eigenvalues of $A$ as $2/3+x_i$, where $\sum_i x_i=0$,
and set $s_2=\sum_i x_i^2$, $p=x_1x_2x_3$. Expansion gives
\begin{equation}\label{t2:eq:volume-expansion}
 \det A'-\det A=\frac{1-g^2}{3}s_2-(1-g^3)p.
\end{equation}
We claim $p\le5s_2/36$. This is immediate if $p\le0$.
If $p>0$, exactly one $x_i$ is positive, so after relabeling
$(x_1,x_2,x_3)=(q+r,-q,-r)$ with $q,r>0$.
The upper eigenvalue bound gives $q+r\le5/6$, whereas
$s_2=2(q^2+qr+r^2)\ge6qr$. Hence
$p=(q+r)qr\le5s_2/36$, as claimed.
Since $g\le1$, \eqref{t2:eq:volume-expansion} now gives
\[
 \det A'-\det A\ge
 \frac{(1-g)s_2}{36}(7+7g-5g^2)\ge0,
\]
because
$7+7g-5g^2=1/9+(g+2/3)(31/3-5g)\ge1/9$ on $[-2/3,1]$.
\end{proof}

\begin{lemma}[Cofactor comparison]\label{t2:lem:cofactor}
If $b,c$ are any two eigenvalues of $A$, and
$b'=\alpha+gb$, $c'=\alpha+gc$, then
\begin{equation}\label{t2:eq:cofactor}
 b'c'\ge f^2bc.
\end{equation}
\end{lemma}
\begin{proof}
We have $0\le b,c\le3/2$ and $b+c\le2$. Also
$\alpha+gx\ge0$ for $0\le x\le3/2$: for $g\ge0$ this is
immediate, and for $g<0$ it follows from
\begin{equation}\label{t2:eq:positive}
 \alpha+gx\ge\alpha+\frac32g=\frac{4+5g}{6}\ge\frac19.
\end{equation}
Write $s=b+c$, $p=bc$.

If $g\ge-1/2$, then $2g+\alpha=2(1+2g)/3\ge0$.
The identity
\[
 4\bigl[b'c'-(g+\alpha)^2bc\bigr]
 =\alpha\bigl[(2g+\alpha)(b-c)^2
       +(2-s)\{2\alpha+(2g+\alpha)s\}\bigr]\ge0
\]
and $(g+\alpha)^2=(2+g)^2/9\ge f^2$ prove the claim.

If $-2/3\le g<-1/2$, Lemma~\ref{t2:lem:multipliers} gives $f^2<1/9$.
Set $h=g^2-1/9>0$. The inequality
$(3/2-b)(3/2-c)\ge0$ gives $p\ge3s/2-9/4$.
Furthermore
\[
 \alpha g+\frac32h=\frac{(5g-1)(g+1)}6<0.
\]
It follows that
\begin{align*}
 b'c'-\frac{bc}{9}
 &=\alpha^2+\alpha gs+hp\\
 &\ge\alpha^2-\frac94h+
       \left(\alpha g+\frac32h\right)s\\
 &\ge\alpha^2+2\alpha g+\frac34h
   =\frac{13+16g-5g^2}{36}\ge\frac{1}{324}>0.
\end{align*}
The last bound follows from
$13+16g-5g^2=1/9+(g+2/3)(58/3-5g)$.
This proves \eqref{t2:eq:cofactor} in both cases.
\end{proof}

We can now show
\begin{equation}\label{t2:eq:det-monotone}
 \det\rho'_2\ge\det\rho_2.
\end{equation}
In a common eigenbasis of $A,A'$, with respective eigenvalues $a,b,c$
and $a',b',c'$, all six eigenvalues are nonnegative. Set
$\Delta=8\det\rho_2$, $\Delta'=8\det\rho'_2$.
Equation~\eqref{t2:eq:cartesian} gives the polynomial identity
\begin{align*}
 abc(\Delta'-\Delta)
 &=(a'b'c'-abc)\Delta\\
 &\quad+aa'r_1^2(b'c'-f^2bc)
       +bb'r_2^2(a'c'-f^2ac)
       +cc'r_3^2(a'b'-f^2ab)\ge0.
\end{align*}
Every term is nonnegative by Lemmas~\ref{t2:lem:volume} and
\ref{t2:lem:cofactor}, and $\Delta\ge0$ follows from $\rho_2\succeq0$.
If $abc>0$, division proves \eqref{t2:eq:det-monotone}.
If $A$ is singular, then
\eqref{t2:eq:cartesian}, $A\succeq0$ and $\det\rho_2\ge0$ force
$\det\rho_2=0$. The inequality is then immediate from
$\rho'_2\succeq0$. This also covers all boundary states.

\clearpage
\section{Determinant comparison at order three}\label{sec:t3}
Let $N\ge4$, $\sigma=\rho_3$ and $\tau=\rho'_3$. We prove
$\det\tau\ge\det\sigma$ by four base cases,
$(N,K)=(4,2),(4,3),(5,4),(6,5)$, and an available-extension convex hull.
Channel vertices, multipliers, auxiliary Kraus maps and spin reversal have
been derived in Section~\ref{sec:prelim}. Every scalar and physical
inequality used for the base cases is proved below. In this chapter
$L=L_4$ is the logarithmic determinant deficit and $V=V_4$ the purity
variance; a separate function $L(s)$ is explicitly defined as a lower
variance envelope when it is introduced.

For later use, on spin $3/2$ set
$\mathcal B(X)=[\operatorname{Tr}(X)I_4-X-\Theta(X)]/2$.
The orthogonality argument in Section~\ref{sec:prelim} proves its
positivity, unitality and trace preservation; its rank multipliers are
$(0,-1,0)$.

\subsection{Scalar determinant bounds}
For a full-rank four-dimensional state put
\[
 L(\omega)=-\log(256\det\omega),\qquad
 V(\omega)=\Tr\omega^2-1/4.
\]
All logarithms are natural.

\begin{lemma}\label{t3:lem:lower}
Every four-dimensional state satisfies
\[
 L(\omega)\ge\frac{49}{8}V(\omega),
\]
with $L=+\infty$ for a singular state.
\end{lemma}
\begin{proof}
On the eigenvalue simplex minimize $F=L-cV$ with $c=49/8$.
It diverges at the boundary, so a minimum lies in the interior.
Stationarity gives $p_i^{-1}+2cp_i=\mathrm{constant}$.
This scalar function is strictly convex and has at most two roots.
Its larger root exceeds $(2c)^{-1/2}=2/7$; if that value occurred
twice, varying those two coordinates oppositely would give the
negative second derivative $2(p_i^{-2}-2c)$.
Thus any nonuniform local minimum has one large entry $a$ and
three equal entries $(1-a)/3$, where $a>1/4$.

Along this family,
\[
 F(a)=-\log\!\left(\frac{256a(1-a)^3}{27}\right)
       -\frac{49}{6}(a-1/4)^2,
\]
and
\[
 F'(a)=(4a-1)\left[\frac1{a(1-a)}-\frac{49}{12}\right].
\]
The two nonuniform stationary points are $3/7$ and $4/7$.
The derivative is positive on $(1/4,3/7)$, negative on
$(3/7,4/7)$ and positive on $(4/7,1)$. The global minimum
therefore occurs either at $a=1/4$, where $F=0$, or at $a=4/7$,
where
\[
 F(4/7)=\log(2401/1024)-27/32>0.
\]
For an exact verification of the final strict inequality, set
$r=(2401-1024)/(2401+1024)=1377/3425$. The positive series
\[
 \log\frac{1+r}{1-r}=2\sum_{k=0}^\infty\frac{r^{2k+1}}{2k+1}
\]
gives
\[
 \log(2401/1024)-\frac{27}{32}
 >2(r+r^3/3)-\frac{27}{32}
 =\frac{4707658629}{1285676500000}>0.
\]
This proves the lemma without numerical optimization.
\end{proof}

\begin{lemma}\label{t3:lem:upper}
If $\lambda_{\min}(\omega)\ge m$, with $0<m<1/4$, then
\[
 L(\omega)\le C(m)V(\omega),\qquad
 C(m)=-\frac{12}{(1-4m)^2}
       \log\!\left(\frac{256m(1-m)^3}{27}\right).
\]
The constant is sharp and is attained at
$(m,(1-m)/3,(1-m)/3,(1-m)/3)$.
\end{lemma}
\begin{proof}
We first prove and use the reciprocal-trace bound
\[
 \Tr\eta^{-1}-16\le\frac{12}{u(1-u)}V(\eta)
       \quad\text{when }\lambda_{\min}(\eta)\ge u.
\]
For completeness, its extremal argument maximizes
$\sum_i p_i^{-1}-16-C\sum_i(p_i-1/4)^2$ on $p_i\ge u$.
Stationarity of free coordinates gives at most two values; the
smaller free value cannot occur twice at a maximum. On the family
with one smaller free value $b$ and $n-1$ equal larger values $a$,
the third derivative is $-6/b^4+6/[(n-1)^2a^4]<0$.
The first derivative is strictly concave as a function of $b$ and
vanishes at the equal-value endpoint. At any earlier zero it is positive
immediately to the right (strict concavity and the later zero), so such
a zero cannot be a local maximum. Thus a nonuniform interior maximum
is impossible. Thus candidates have
$k$ entries $u$ and $4-k$ equal entries $(1-ku)/(4-k)$.
Their ratio $(\Tr\eta^{-1}-16)/V(\eta)$ is
$4(4-k)/[u(1-ku)]$, decreasing in $k$, and is largest at $k=1$.

Apply this bound to
$\eta_s=(\omega+sI)/(1+4s)$, whose minimum eigenvalue is at least
$u_s=(m+s)/(1+4s)$. It gives
\[
 \Tr(\omega+sI)^{-1}-\frac4{s+1/4}
 \le\frac{12V(\omega)}{(1+4s)(m+s)(1-m+3s)}.
\]
Integrating over $s\in[0,\infty)$ yields $L(\omega)$ on the left.
The upper reciprocal bound is attained simultaneously at every
$s$ by the spectrum $(m,(1-m)/3,(1-m)/3,(1-m)/3)$.
Hence the integral on the right equals that spectrum's ratio
$L/V$, namely the advertised $C(m)$. This also proves sharpness.
\end{proof}

\paragraph{A sufficient determinant criterion.}
Let $\sigma,\tau$ be physical four-dimensional states with
$V(\tau)\le q^2V(\sigma)$ and $\lambda_{\min}(\tau)\ge m>0$.
If
\[
 q^2C(m)\le49/8,
\]
then $\det\tau\ge\det\sigma$. For positive-definite $\sigma$,
\[
 L(\tau)\le C(m)V(\tau)\le q^2C(m)V(\sigma)
         \le(49/8)V(\sigma)\le L(\sigma).
\]
A singular input has determinant zero, so its case is immediate.


\subsection{Positive output floors for three base channels}
For $N=4,5,6$, define $G_q^{(N)}$, $0\le q\le3$, to be the
Clebsch--Gordan channel from spin $j=N/2$ to spin $a=3/2$ with
environment spin $b=(N-3)/2+q$. Explicitly,
\[
 G_q^{(N)}(\rho)=\sum_{\mu=-b}^b A_\mu\rho A_\mu^*,\qquad
 (A_\mu)_{m,n}=\langle a,m;b,\mu\mid j,n\rangle.
\]
The Clebsch--Gordan isometry proves complete positivity and trace
preservation. Covariance gives $G_q^{(N)}(I_{N+1}/(N+1))=I_4/4$.
Here $G_0^{(N)}=R_{N,3}$. Relative to $R_{N,3}$, the multipliers of
$G_1,G_2$ are
\[
\begin{array}{c|ccc|ccc}
 N&\multicolumn{3}{c|}{G_1}&\multicolumn{3}{c}{G_2}\\\hline
4&2/3&0&-1&1/9&-5/7&3/7\\
5&13/21&-1/7&-9/7&1/21&-11/14&9/14\\
6&7/12&-1/4&-3/2&0&-5/6&5/6
\end{array}
\]
These entries follow directly from the displayed finite CG sums;
\path{verify_positive_floors.py} checks them with exact algebraic
coefficients. Let $D(\rho)=\Tr(\rho)I_4/4$, and let
$\mathcal B(X)=[\Tr(X)I_4-X-\Theta X]/2$ be the positive map
described below. Exact comparison of the three multipliers gives
\begin{align}
 R_{4,3}\Phi^{(4)}_3
 &=\frac{11}{70}D+\frac17\mathcal B G_0^{(4)}
             +\frac1{10}\Theta G_1^{(4)}+\frac35G_2^{(4)},\\
 R_{5,3}\Phi^{(5)}_4
 &=\frac{46}{189}D+\frac{170}{567}\mathcal B G_0^{(5)}
             +\frac{140}{567}\Theta G_1^{(5)}+\frac{119}{567}G_2^{(5)},\\
 R_{6,3}\Phi^{(6)}_5
 &=\frac7{30}D+\frac1{25}\Theta G_0^{(6)}
             +\frac{11}{30}\mathcal B G_0^{(6)}+\frac9{25}\Theta G_1^{(6)}.
\end{align}
Every summand maps positive inputs to positive outputs, and the
coefficients in each line sum to one. Consequently the respective
output lower bounds are
\[
 \tau_{4,3}\ge\frac{11}{280}I_4,\qquad
 \tau_{5,4}\ge\frac{23}{378}I_4,\qquad
 \tau_{6,5}\ge\frac7{120}I_4.
\]
Together with the scalar log-determinant bounds, these prove the
three determinant inequalities for arbitrary mixed inputs.


The preceding output floors and exact multiplier
contractions are
\[
\begin{array}{c|c|c}
 (N,K)&m& q^2\\\hline
 (4,3)&11/280&16/49\\
 (5,4)&23/378&1/4\\
 (6,5)&7/120&1/4
\end{array}
\]
The floor certificates express a radial expansion of the induced
map as a convex combination of physical $N\to3$ covariant channels
and their compositions with positive local maps. Their exact positive decompositions were displayed above.

\paragraph{The $(4,3)$ case.}
For $m=11/280$,
\[
 \frac{16}{49}C(m)
 =\frac{19200}{3481}
     \log\frac{648270000}{214116199}.
\]
The logarithm is strictly below $111/100$. Here is a rational
certificate. First,
\[
 \frac{109}{36}-\frac{648270000}{214116199}
 =\frac{945691}{7708183164}>0.
\]
Set $r=73/145$, so $(1+r)/(1-r)=109/36$.
Bounding the positive logarithmic-series tail geometrically gives
\[
 \log\frac{109}{36}
 \le U:=2(r+r^3/3)+\frac{2r^5}{5(1-r^2)},
 \qquad \frac{111}{100}-U
 =\frac{84307637}{119628045000}>0.
\]
Consequently
\[
 \frac{16}{49}C(11/280)
 <\frac{21312}{3481}<\frac{49}{8},\qquad
 \frac{49}{8}-\frac{21312}{3481}=\frac{73}{27848}>0.
\]
The sufficient criterion proves determinant monotonicity for every
input in this case.

\paragraph{The $(5,4)$ and $(6,5)$ cases.}
Both output floors exceed $1/20$. Since the admissible spectral
set shrinks when the minimum-eigenvalue bound increases, the sharp
constant $C(m)$ is nonincreasing. Moreover
\[
 C(1/20)=\frac{75}{4}\log\frac{16875}{6859}<18.
\]
An exact certificate is
$\sum_{k=0}^3(24/25)^k/k!>16875/6859$, which implies
$\log(16875/6859)<24/25$.
Thus $q^2C(m)<18/4=9/2<49/8$ in both cases, proving their
determinant inequalities for every input state.


\subsection{Physical spin-two inputs: Clifford reduction for the K=2 vertex}
Let $\rho$ be a spin-two density matrix and let
\[
 \sigma=\rho_3=\frac{I_4}{4}+B_1+Q+B_3
\]
be its three-body marginal, decomposed into tensor ranks one, two, and three.
Set
\[
 v=\operatorname{Tr}B_1^2,\quad u=\operatorname{Tr}Q^2,\quad
 w=\operatorname{Tr}B_3^2,\quad r=\sqrt{5v},\quad z=\sqrt u,
 \quad s=r+3z.
\]
Thus $r=|\langle\mathbf J_{3/2}\rangle_\sigma|$ and
$V=\operatorname{Tr}\sigma^2-1/4=r^2/5+z^2+w$.

\paragraph{Two positive matrices supplied by the physical extension.}
Besides $\sigma\ge0$, the spin-two to spin-$3/2$ extremal channel with
spin-$3/2$ environment gives the positive matrix
\[
 Y=\frac{I_4}{4}+C\ge0,\qquad C=\frac23B_1-B_3.
\]
Its induced multiplier triple is $(2/3,0,-1)$; this follows directly from
Clebsch--Gordan recoupling. This channel is applied to the original spin-two
state, so its positivity on the physical marginal domain is automatic.
No positivity on arbitrary spin-$3/2$ states is assumed.

Spin-$3/2$ time reversal $\Theta(X)=VX^{\mathsf T}V^\dagger$ preserves
rank-two tensors and negates ranks one and three. Therefore
$\Theta(Y)=I_4/4-C\ge0$, giving
\begin{equation}\label{t3:eq:C-constraint}
 \|C\|_{\rm op}\le\frac14,\qquad
 \frac49v+w=\operatorname{Tr}C^2\le\frac14.
\end{equation}
Also $[\sigma+\Theta(\sigma)]/2=I_4/4+Q\ge0$.
The five spin-$3/2$ quadrupoles form a Clifford vector: if $\Gamma_i$ are
normalized by $\{\Gamma_i,\Gamma_j\}=2\delta_{ij}I_4$, then
$Q=\frac12\sum_i q_i\Gamma_i$, $u=\sum_iq_i^2$, and
\[
 Q^2=\frac u4I_4,\qquad\|Q\|_{\rm op}=\frac z2,\qquad0\le z\le\frac12.
\]
One explicit choice of the five $\Gamma_i$ is
$\{J_y,J_z\}/\sqrt3$, $\{J_z,J_x\}/\sqrt3$,
$\{J_x,J_y\}/\sqrt3$, $(J_x^2-J_y^2)/\sqrt3$, and
$J_z^2-5I_4/4$.

\paragraph{Additional elementary constraints from $\sigma\ge0$.}
Rotate the mean spin onto the positive $z$ axis. Then $B_1=(r/5)J_z$,
$0\le r\le3/2$, and $\|B_1\|_{\rm op}=3r/10$.
With the Hilbert--Schmidt normalized diagonal tensors
\[
 T_{20}=\frac12\operatorname{diag}(1,-1,-1,1),\qquad
 T_{30}=\frac1{2\sqrt5}\operatorname{diag}(1,-3,3,-1),
\]
the elementary diagonal operator inequalities
\[
 J_z-T_{20}\le I_4,\qquad
 J_z-\frac{\sqrt5}{2}T_{30}\le\frac54I_4
\]
yield
\begin{equation}\label{t3:eq:rz-constraints}
 r\le1+z,\qquad
 w\ge\frac{(2r-5/2)_+^2}{5}.
\end{equation}
Here $x_+=\max\{x,0\}$. These follow from
$|\operatorname{Tr}(\sigma T_{20})|\le z$ and
$|\operatorname{Tr}(\sigma T_{30})|\le\sqrt w$.
Together, $0\le r\le3/2$ and $0\le z\le1/2$ imply $0\le s\le3$.


\paragraph{An input-dependent lower eigenvalue bound.}
For $N=4,K=2$, the target marginal is
\begin{align*}
 \tau&=\frac{I_4}{4}+\frac12B_1-\frac3{14}Q-\frac47B_3\\
 &=\frac{I_4}{4}-\frac3{14}Q+\frac47C+\frac5{42}B_1.
\end{align*}
The operator norms above give the exact bound
\begin{equation}\label{t3:eq:K42-output-floor}
 \lambda_{\min}(\tau)
 \ge\frac14-\frac{3z}{28}-\frac17-\frac r{28}
 =\frac{3-s}{28}.
\end{equation}
For $s<3$ this is strictly positive. At $s=3$, one necessarily has
$r=3/2$ and $z=1/2$; maximal spin expectation forces $\sigma$ to be a
pure coherent state, so $\det\sigma=0$ and the unshifted determinant
comparison is automatic from $\tau\ge0$.

\paragraph{Reduction to one real parameter.}
The input variance obeys $V\ge L(s)$, where
\begin{equation}\label{t3:eq:input-variance-lower}
 L(s)=\begin{cases}
 s^2/14,&0\le s\le7/3,\\[2pt]
 \displaystyle\frac14+\frac{(s-3/2)^2}{5},&7/3\le s\le11/4,\\[5pt]
 \displaystyle\frac14+\frac{(s-3/2)^2+(2s-11/2)^2}{5},&11/4\le s\le3.
 \end{cases}
\end{equation}
For the first interval, this is Cauchy--Schwarz:
$(r+3z)^2\le14(r^2/5+z^2)$. Its unconstrained minimizer
$r=5s/14,z=3s/14$ reaches $z=1/2$ at $s=7/3$.
Beyond that point minimize
\[
 \frac{r^2}{5}+\frac{(s-r)^2}{9}
       +\frac{(2r-5/2)_+^2}{5}
\]
over $r\ge s-3/2$. Its derivative is nonnegative there for $s\ge7/3$,
so the minimum occurs at $r=s-3/2,z=1/2$. The final term becomes active
at $s=11/4$, proving the displayed formula.

The output variance is
\[
 V_\tau=\operatorname{Tr}\tau^2-1/4
 =\frac{r^2}{20}+\frac9{196}z^2+\frac{16}{49}w.
\]
Using~\eqref{t3:eq:C-constraint} gives
\[
 V_\tau\le G_s(r):=\frac4{49}+\frac{37r^2}{1764}
                        +\frac{(s-r)^2}{196}.
\]
This is convex in $r$. On the enlarged interval
$\max\{0,s-3/2\}\le r\le\min\{s,3/2\}$, its larger endpoint value is
\begin{equation}\label{t3:eq:output-variance-upper}
 U(s)=\begin{cases}
 \displaystyle\frac4{49}+\frac{37s^2}{1764},&0\le s\le3/2,\\[5pt]
 \displaystyle\frac{101+4(s-3/2)^2}{784},&3/2\le s\le3.
 \end{cases}
\end{equation}
In the second interval the endpoint $r=3/2$ dominates $r=s-3/2$, since
the difference is $1/28-(s-3/2)^2/63\ge0$.
Equivalently, with $\varepsilon=3-s$,
\[
 U(s)=\frac{55-6\varepsilon+2\varepsilon^2}{392}
 \qquad(3/2\le s\le3).
\]
In particular, $V_\tau\le55/392$ and
$\operatorname{Tr}\tau^2\le153/392$, with equality at a coherent input.
The bounds~\eqref{t3:eq:K42-output-floor}, \eqref{t3:eq:input-variance-lower}, and
\eqref{t3:eq:output-variance-upper} are rigorous consequences of the two
positive matrices $\sigma$ and $Y$, plus the spin-$3/2$ Clifford identities.

\subsubsection{The region $0\le s\le13/5$}
Put $s=r+3z$ and $h=3z$, so $0\le r,h\le3/2$.
The output minimum eigenvalue is at least $m(s)=(3-s)/28$.
For $V>0$, maximizing the purity ratio over the allowed octupole
norm $w$ gives
\[
 \frac{V(\tau)}{V(\sigma)}
 \le Q(r,h):=\frac{144+37r^2+9h^2}{49(9+4r^2+4h^2)}.
\]
Indeed the ratio increases with $w$, whose upper bound is
$1/4-4r^2/45$. Independently the ratio is at most $16/49$.
The following rational bounds suffice:
\[
\begin{array}{c|c|c|c}
 s\text{ interval}&V(\tau)/V(\sigma)\le b&C(m(s))<c&bc\\\hline
 {[0,8/5]}&16/49&18&288/49\\
 {[8/5,2]}&4/15&20&16/3\\
 {[2,12/5]}&1/4&24&6\\
 {[12/5,5/2]}&23/100&26&299/50\\
 {[5/2,13/5]}&11/50&277/10&3047/500
\end{array}
\]
Every entry in the last column is strictly less than $49/8$.
It follows that
\[
 L(\tau)\le C(m(s))V(\tau)<\frac{49}{8}V(\sigma)\le L(\sigma)
\]
whenever $V(\sigma)>0$; the maximally mixed case gives equality.

Here are reproducible exact certificates for the table. Along each
ray $(r,h)\mapsto(k r,k h)$, $Q$ decreases with $k>0$. Thus to
bound $Q$ for $r+h\ge s_0$, it suffices to consider $r+h=s_0$.
The inequality $Q\le b$ is the quadratic inequality
\[
 441b-144+(196b-37)r^2+(196b-9)(s_0-r)^2\ge0
\]
on $\max(0,s_0-3/2)\le r\le\min(3/2,s_0)$.
Its minimum is obtained by clamping the rational critical point
$(196b-9)s_0/(392b-46)$ to this interval. Substitution gives a
strictly positive rational number in each of the last four rows.
For the $C$ column, monotonicity allows substitution of the upper
endpoint of the $s$ interval. Each resulting logarithm is bounded
using the rational inequality
\[
 \log x\le 2\sum_{k=0}^{M-1}\frac{q^{2k+1}}{2k+1}
          +\frac{2q^{2M+1}}{(2M+1)(1-q^2)},
 \qquad q=\frac{x-1}{x+1},\quad M=40.
\]
The right side is an upper bound because the omitted positive
series is bounded by a geometric tail. The embedded \path{determinant_t3/bounds/verify_central_region.py} checks every rational inequality
and records the exact quadratic minima and margins in
\path{central_region_certificates.json}.

\subsubsection{The region $13/5\le s<3$}
Put $\epsilon=3-s\in(0,2/5]$. The preceding reduction gives
\begin{align}
 V(\tau)&\le U(\epsilon)=\frac{55-6\epsilon+2\epsilon^2}{392},\label{t3:eq:out}\\
 V(\sigma)&\ge v_*(\epsilon)=
 \begin{cases}3/4-\epsilon+\epsilon^2,&0\le\epsilon\le1/4,\\
 1/4+(3/2-\epsilon)^2/5,&1/4\le\epsilon\le2/5.
 \end{cases}\label{t3:eq:in}
\end{align}
\begin{lemma}[Exact minimum determinant deficit at fixed purity]
For a four-dimensional state with $V(\sigma)=v\in(0,3/4)$,
\[
 L(\sigma):=-\log(256\det\sigma)\ge
 \ell(v):=-\log\!\left[256a\left(\frac{1-a}{3}\right)^3\right],
 \quad a=\frac14+\sqrt{\frac{3v}{4}}.
\]
The right side increases with $v$.
\end{lemma}
\begin{proof}
Minimize $L$ under fixed trace and purity. A minimum is interior,
because $L$ diverges on a singular spectrum while a full-rank
spectrum with the prescribed purity exists. Lagrange multipliers
give at most two eigenvalue values, say $a>b$.
The Hessian of the Lagrangian at a coordinate equal to $a$ is
$1/a^2-1/(ab)<0$. If $a$ occurred twice, an opposite variation
of those two coordinates is tangent to both constraints and gives
a negative second variation, excluding a minimum. Hence the
minimizing spectrum has one large value $a$ and three equal values
$(1-a)/3$. The purity constraint gives the displayed $a$.
The derivative of $-\log[256a(1-a)^3/27]$ is
$(4a-1)/[a(1-a)]>0$ for $a>1/4$, proving monotonicity.
\end{proof}

The sharp upper Burg bound established in Lemma~\ref{t3:lem:upper} gives
\[
 L(\tau)\le C(\epsilon/28)U(\epsilon),\qquad
 C(m)=-\frac{12}{(1-4m)^2}\log\frac{256m(1-m)^3}{27}.
\]
It therefore suffices to prove the one-variable inequality
\begin{equation}\label{t3:eq:goal}
 C(\epsilon/28)U(\epsilon)<\ell(v_*(\epsilon))
                  \quad(0<\epsilon\le2/5).
\end{equation}

\subsubsection{The endpoint region $0<\epsilon\le1/8$}
In this regime, \eqref{t3:eq:in} implies
$a\ge1-\epsilon/2$, because
\[
 \frac34v_*(\epsilon)
 =\frac9{16}-\frac{3\epsilon}{4}+\frac{3\epsilon^2}{4}
 \ge(3/4-\epsilon/2)^2.
\]
Since the determinant of the one-large-eigenvalue spectrum decreases
with $a>1/4$,
\[
 L(\sigma)\ge
 -\log\!\left[\frac{32}{27}(1-\epsilon/2)\epsilon^3\right]
 >3\log(1/\epsilon)-\frac15.
\]
Here $\log(32/27)<1/5$ follows from
$\exp(1/5)>1+1/5>32/27$.

On the other hand, $U(\epsilon)\le55/392$, and
\begin{align*}
 C(\epsilon/28)U(\epsilon)
 &\le\frac{12(55/392)}{(1-\epsilon/7)^2}
 \left[\log(1/\epsilon)+\log\frac{189}{64}
                         -3\log(1-\epsilon/28)\right]\\
 &< A\left[\log(1/\epsilon)+\frac{11}{10}\right],
       \qquad A=\frac{1056}{605}.
\end{align*}
Indeed $(1-\epsilon/7)^{-2}\le(56/55)^2$, and
\[
 \log\!\left[\frac{189}{64}\left(\frac{224}{223}\right)^3\right]
 <\frac{11}{10}.
\]
The latter inequality is certified rationally by
$\sum_{k=0}^6(11/10)^k/k!>
 (189/64)(224/223)^3$.
Finally $\log(1/\epsilon)\ge\log8>2$, since
$\log2=2\sum_{k\ge0}(1/3)^{2k+1}/(2k+1)>2/3$.
Consequently
\[
 L(\sigma)-L(\tau)
 >2(3-A)-\frac{11A}{10}-\frac15
 =\frac{1177}{3025}>0.
\]
At $\epsilon=0$, \eqref{t3:eq:in} forces $V(\sigma)=3/4$;
the input is pure and its determinant is zero, so the desired
determinant inequality is immediate.

\subsubsection{The remaining interval $1/8\le\epsilon\le2/5$}
Both $C(\epsilon/28)$ and $U(\epsilon)$ are positive and
nonincreasing. The first monotonicity follows from the sharp
variational characterization of $C(m)$: increasing the minimum
eigenvalue floor shrinks the maximization domain.
Also $v_*(\epsilon)$ decreases with $\epsilon$, while $\ell$ increases.
For an interval $[a,b]$ it therefore suffices to verify
\[
 C(a/28)U(a)<\ell(v_*(b)).
\]
The script \path{certify_k42_outer.py} verifies this strict
inequality on the 275 intervals
\[
 [k/1000,(k+1)/1000],\qquad k=125,126,\ldots,399,
\]
using rational arithmetic only. The smallest certified difference
exceeds $0.02341839$ and occurs on $[399/1000,2/5]$.
That decimal is descriptive; the comparisons themselves are exact.

Square roots are bounded below with denominator $2^{40}$ by an
integer square root. For $q>0$, set
\[
 h=2^{-40}\left\lfloor\sqrt{2^{80}q}\right\rfloor;
 \qquad h^2\le q<(h+2^{-40})^2.
\]
Using $a_0=1/4+h\le1/4+\sqrt{3v_*(b)/4}$ gives a lower bound
on the input logarithmic deficit.

To bound a logarithm, first write its positive rational argument as
$2^kR$ with $1\le R<2$. For $z=(R-1)/(R+1)\in[0,1/3)$,
\[
 2\sum_{j=0}^{11}\frac{z^{2j+1}}{2j+1}
 \le\log R\le
 2\sum_{j=0}^{11}\frac{z^{2j+1}}{2j+1}
   +\frac{2z^{25}}{25(1-z^2)}.
\]
The same bounds with $z=1/3$ bound $\log2$. Thus both endpoints
of every logarithmic enclosure are explicit rational numbers.
The script saves each interval, the rational square-root bound,
and the strictly positive rational gap in
\path{k42_outer_certificates.json}.

This completes \eqref{t3:eq:goal} and hence the $(4,2)$ determinant
inequality throughout $13/5\le r+3z\le3$. Together with the
independent central-regime proof, it establishes the final finite
base case required by the all-spin order-three reduction.
\subsection{The exact reduction to four finite-spin cases}
The four base cases have now been proved. We finish by establishing
the exact reduction from arbitrary spin to those cases.

For $N\ge4$, let $\Phi_K^{(N)}$ be the extremal covariant channel on
spin $N/2$, and write $R_{N,3}$ for reduction to three symmetric qubits.
Let $v_{N,K}=(f_1,f_2,f_3)$ be the induced multiplier triple.

\begin{theorem}\label{t3:thm:reduction}
The following statements are equivalent.
\begin{enumerate}
\item For every $N\ge4$, every spin-$N/2$ state $\rho$, and every
covariant channel $\Phi$ on that sector,
\[
 \det(R_{N,3}\Phi(\rho))\ge\det(R_{N,3}\rho).
\]
\item The same determinant inequality holds for all input states in
each of the following four cases:
\[
 (N,K)=(4,2),\qquad(4,3),\qquad(5,4),\qquad(6,5).
\]
\end{enumerate}
\end{theorem}

The sufficiency uses exact rational certificates on a continuous
parameter domain covering all $N\ge7$, plus exact finite checks for
$N=4,5,6$.

\subsubsection{The positive-map polytope on spin $3/2$}
Here ``positive map'' means positivity, not complete positivity.
Spin reversal is the linear map
\[
 \Theta(X)=e^{-i\pi J_y}X^{\mathsf T}e^{i\pi J_y}.
\]
It is positive, unital, trace preserving and covariant; its rank-$\ell$
multiplier is $(-1)^\ell$. On a pure state $P_\psi$, the vectors
$\psi$ and its time-reversed vector are orthogonal, since the spin
is half-integral. Therefore
\[
 \mathcal B(X)=\tfrac12[\Tr(X)I_4-X-\Theta(X)]
\]
is also positive, unital, trace preserving and covariant:
$\mathcal B(P_\psi)=(I_4-P_\psi-\Theta(P_\psi))/2\ge0$.
Its multiplier triple is $(0,-1,0)$.

Let $\Psi_1$ denote the local spin-$3/2$ rank-one covariant vertex.
The five positive maps
\[
 \mathrm{id},\quad\Theta,\quad\Psi_1,\quad\Theta\Psi_1,\quad\mathcal B
\]
have multiplier triples
\begin{equation}\label{t3:eq:positivevertices}
 (1,1,1),\quad(-1,1,-1),\quad
 (11/15,1/5,-3/5),\quad(-11/15,1/5,3/5),\quad(0,-1,0).
\end{equation}
Their convex hull is exactly the set of positive, unital,
trace-preserving covariant maps on spin $3/2$. Its six facets are
\begin{align}
 |9f_1+f_3|&\le5(1+f_2),\label{t3:eq:coherent1}\\
 3|f_1-f_3|&\le5(1-f_2),\label{t3:eq:coherent2}\\
 2|f_3|&\le1+f_2.\label{t3:eq:ghz}
\end{align}
Sufficiency follows from the exact facet/vertex description and the
positivity of the five generating maps. Necessity follows from two
test states. A coherent input has output eigenvalues
\[
 \frac{5+9f_1+5f_2+f_3}{20},\quad
 \frac{5+3f_1-5f_2-3f_3}{20},\quad
 \frac{5-3f_1-5f_2+3f_3}{20},\quad
 \frac{5-9f_1+5f_2-f_3}{20},
\]
giving \eqref{t3:eq:coherent1}--\eqref{t3:eq:coherent2}. The spin-$3/2$ GHZ
state has output spectrum
\[
 \left(\frac{1+f_2+2f_3}{4},\frac{1+f_2-2f_3}{4},
               \frac{1-f_2}{4},\frac{1-f_2}{4}\right),
\]
giving \eqref{t3:eq:ghz}. The additional inequality $f_2\le1$ is
already implied by \eqref{t3:eq:coherent2}.

For physical induced maps, the four coherent inequalities hold
automatically. Thus such an induced map is positive if and only if
$2|f_3|\le1+f_2$.

\subsubsection{A thirteen-vertex hull containing every physical triple}
Let $\mathcal H$ be the convex hull of the five triples
\eqref{t3:eq:positivevertices}, the four triples
\begin{align*}
 v_{4,2}&=(1/2,-3/14,-4/7),\\
 v_{4,3}&=(0,-4/7,5/14),\\
 v_{5,4}&=(-1/7,-1/2,19/42),\\
 v_{6,5}&=(-1/4,-5/12,1/2),
\end{align*}
and their four spin-reversed triples, obtained by reversing the signs
of $f_1,f_3$. Its twenty facets can be written as ten pairs:
\begin{align*}
 |61f_1+174f_3|&\le110f_2+125,&
 |13f_1+129f_3|&\le57f_2+85,\\
 |9f_1+f_3|&\le5f_2+5,&
 |5f_1-8f_3|&\le9f_2+9,\\
 |4f_1+6f_3|&\le5f_2+5,&
 3|f_1-f_3|&\le5-5f_2,\\
 |3f_1-19f_3|&\le3f_2+13,&
 |3f_1+12f_3|&\le10f_2+10,\\
 |f_1-3f_3|&\le3f_2+3,&
 17|f_3|&\le6f_2+11.
\end{align*}
The script \path{determinant_t3/search/audit_reduction.py} verifies the equivalence exactly:
it enumerates all intersections of triples of facet planes, retains
the feasible intersections and obtains precisely the thirteen
listed vertices. The four coherent facets make the polyhedron
bounded: a recession vector $h$ must obey
$5h_2\ge|9h_1+h_3|$ and $-5h_2\ge3|h_1-h_3|$, hence $h=0$.
Thus no numerical convex-hull assertion is needed for the certificate.

\paragraph{All $N\ge7$.}
Set $D=N(N+2)$, $u=K(K+1)$, $y=1/(N+2)$ and $x=f_1=1-2u/D$.
For every physical vertex, $0<y\le1/9$ and $-1+2y\le x\le1$.
The other multipliers are
\[
 f_2=\frac{3Dx^2-D-6x}{2(D-3)},\qquad
 f_3=\frac{5D^2x^3-3D^2x-40Dx^2+8Dx+10D+48x}
                  {2(D-8)(D-3)},\qquad D=\frac{1-2y}{y^2}.
\]
The four coherent facets hold automatically. For each of the other
sixteen facets, multiply its nonnegative-side expression by
\[
 Q(y)=2(1-3y)(1+y)(1-4y)(1+2y)>0
           \qquad(0\le y\le1/9)
\]
and substitute
\[
 y=Y/9,\qquad x=-1+2y+2(1-y)Z,
             \qquad (Y,Z)\in[0,1]^2.
\]
The resulting polynomial is nonnegative on the whole unit square.
This is certified by exact nonnegative Bernstein coefficients on
the dyadic boxes saved in \path{finite_hull_certificates.json}.
There are sixteen certificates, using between four and twenty-two
boxes each. All coefficients and endpoints are rational.

For a polynomial $P(Y,Z)=\sum a_{ij}Y^iZ^j$ of bidegree $(n,m)$,
its Bernstein coefficient indexed by $(k,l)$ is
\[
 b_{kl}=\sum_{i\le k,j\le l}a_{ij}
          \frac{\binom{k}{i}}{\binom{n}{i}}
          \frac{\binom{l}{j}}{\binom{m}{j}}.
\]
Nonnegative coefficients prove nonnegativity because the Bernstein
basis functions are nonnegative and sum to one. The script
\path{determinant_t3/search/audit_reduction.py} carries out this computation after
each exact affine box substitution and verifies the entire partition
by recursive subdivision. This proves $v_{N,K}\in\mathcal H$ for
every $N\ge7$ and every $K$, not merely for finitely many sampled spins.

\paragraph{Small spins and the availability of extensions.}
All $N=6$ vertices lie in $\mathcal H$, as checked exactly. At $N=4$,
the only nonpositive induced vertices are the two designated base
cases. At $N=5$, the only additional case apart from $(5,4)$ is
$(5,3)$, and it uses only maps available at lower spins:
\[
 v_{5,3}=\frac1{63}(1,1,1)+\frac1{42}(-1,1,-1)
          +\frac{199}{630}(0,-1,0)+\frac{29}{45}v_{4,2}.
\]
All four coefficients are nonnegative and sum to one. These finite
statements are also checked symbolically by \path{determinant_t3/search/audit_reduction.py}.

\subsubsection{Why the hull proves the finite reduction}
Assume the four determinant inequalities in Theorem~\ref{t3:thm:reduction}.
For a spin-$N/2$ input with $N\ge6$, let $\sigma=R_{N,3}\rho$.
For every base spin $m\in\{4,5,6\}$, the state
$\rho_m=R_{N,m}\rho$ is a physical spin-$m/2$ state with the same
three-qubit marginal $\sigma$. Therefore applying a base induced map
to $\sigma$ gives a positive output and cannot lower its determinant,
by the assumed base inequality. A spin-reversed base output has the
same spectrum and determinant. Each of the five positive local maps
also cannot lower the determinant, because a positive unital
trace-preserving map majorizes spectra.

Any target triple in $\mathcal H$ is a convex combination of these
thirteen maps. Its output on $\sigma$ is the corresponding convex
combination of their outputs. If $\det\sigma>0$, every component
output is positive definite and log-determinant concavity gives
\[
 \log\det\Bigl(\sum_iw_i\tau_i\Bigr)
 \ge\sum_iw_i\log\det\tau_i\ge\log\det\sigma.
\]
If $\det\sigma=0$, the determinant inequality is immediate from
positivity of the target output. The same argument extends from
full-channel vertices to arbitrary convex mixtures of vertices.
For $N=4,5$, use the small-spin reductions above, which avoid any
unavailable higher-spin extension. This proves sufficiency.
Necessity is immediate since the four cases are instances of the
all-spin statement.

\clearpage
\section{Determinant comparison at order four}\label{sec:t4}
Let $N\ge5$, $\sigma=\rho_4$ and $\tau=\rho'_4$. We prove
$\det\tau\ge\det\sigma$ by fourteen base cases and an available-extension
convex hull. Thirteen base cases use the same physical spectral-floor
method, with two analytic endpoint arguments. The fourteenth,
$(N,K)=(5,2)$, uses an exact positive weighted polynomial identity.
All ingredients of both methods are supplied below and in the embedded
data appendix. The exact scalar partition and the positive identity are checked
independently of the Lean certification.

\subsection{The exact all-spin finite reduction}
We first prove the fourteen-case reduction, then the thirteen spectral-floor
cases, and finally the exceptional case by an exact weighted sum of squares.

Let $R_{N,4}$ be reduction from spin $N/2$ to four symmetric qubits,
and let $\Phi_K^{(N)}$, $K=0,\ldots,N$, denote the extremal covariant
channels on spin $N/2$. The proper-order domain is $N\ge5$.

\begin{theorem}
The inequality
\[
 \det(R_{N,4}\Phi(\rho))\ge\det(R_{N,4}\rho)
\]
holds for every $N\ge5$, every state $\rho$ and every covariant
channel $\Phi$ if and only if it holds for every input state in
the following fourteen extremal cases:
\[
\begin{array}{c|l}
 N&K\\\hline
 5&2,3,4\\
 6&3,4,5\\
 7&4,5,6\\
 8&4,7\\
 9&8\\
 10&9\\
 11&10
\end{array}
\]
\end{theorem}

The sufficiency is certified exactly for all spins, through
integer convex-hull calculations, rational Bernstein certificates
on a continuous domain, and exact finite decompositions respecting
the available marginal extensions.

\subsubsection{Multipliers and the finite hull}
Write $v_{N,K}=(f_1,f_2,f_3,f_4)$ for the induced multiplier vector.
For $D=N(N+2)$ and $u=K(K+1)$,
\begin{align*}
 f_1={}&1-2u/D,\\
 f_2={}&1-6u/D+6u(u-2)/[D(D-3)],\\
 f_3={}&1-12u/D+30u(u-2)/[D(D-3)]
               -20u(u-2)(u-6)/[D(D-3)(D-8)],\\
 f_4={}&1-20u/D+90u(u-2)/[D(D-3)]
               -140u(u-2)(u-6)/[D(D-3)(D-8)]\\
       &+70u(u-2)(u-6)(u-12)/[D(D-3)(D-8)(D-15)].
\end{align*}
These are the terminating Racah multipliers.
Partial trace intertwines tensor ranks, so they also specify the
induced map on the spin-two marginal.

Spin reversal $\Theta(X)=e^{-i\pi J_y}X^{\mathsf T}e^{i\pi J_y}$
is positive, unital and trace preserving. It acts on multiplier
vectors by
\[
 (f_1,f_2,f_3,f_4)\mapsto(-f_1,f_2,-f_3,f_4).
\]
Let $\mathcal H$ be the convex hull of the local spin-two covariant
vertices $v_{4,K}$, their spin reversals, and all $v_{N,K}$ with
$5\le N\le11$ and their spin reversals.
Its exact extreme set consists of 34 vectors:
\begin{itemize}
\item the eight local positive vectors $v_{4,K}$ and their reversals
      for $K=0,1,2,3$;
\item the thirteen physical vectors in the theorem's table except
      $(8,4)$, and their thirteen reversals.
\end{itemize}
The local $K=4$ vector is redundant.
The full hull has exactly 161 facets, listed as primitive integer
coefficient vectors in \path{exact_hull.json}.

\paragraph{Exact and complete facet verification.}
Floating-point geometry is used only to propose a retained vertex
list. Every omitted input vector is then verified as an exact
rational convex combination of the retained list; there are 88
such decompositions. Thus the proposed list provably generates the
same hull. The retained points affinely span dimension four.

The script next enumerates \emph{every} quadruple of the 34 retained
points, using integer homogeneous coordinates. Cofactors give the
exact candidate supporting hyperplane; a candidate is retained only
when all vertices lie on one side. All $\binom{34}{4}=46376$
quadruples are tested. Every facet of a full-dimensional four-dimensional
polytope contains four affinely independent vertices, so this process
produces the complete facet list. No floating-point assertion about
facet completeness is used in the proof.

\subsubsection{All large spins lie in this hull}
First consider $N\ge12$ and $K\ge4$. Introduce
\[
 y=\frac1{N+2},\qquad x=f_1=1-\frac{2K(K+1)}{N(N+2)},
 \qquad D=\frac{1-2y}{y^2}.
\]
The physical range obeys
\[
 0<y\le1/14,\qquad
 -1+2y\le x\le1-\frac{40y^2}{1-2y}.
\]
The restriction $K\ge4$ is used here. Extending the parameter range
all the way to $x=1$ would include nonphysical intermediate values
between the low-$K$ vertices and would invalidate some hull inequalities.
The low-$K$ cases are handled separately below.

Five facets are automatic from positivity of the output on a global
coherent input. In integer coefficient form, their nonnegative-side
expressions are
\begin{align*}
 14+28f_1+20f_2+7f_3+f_4,&\qquad
 7+7f_1-5f_2-7f_3-2f_4,\\
 7-10f_2+3f_4,&\qquad
 7-7f_1-5f_2+7f_3-2f_4,\\
 14-28f_1+20f_2-7f_3+f_4.&
\end{align*}
These are positive multiples of the five output probabilities.

For the other 156 facets, substitute
\[
 y=Y/14,\qquad
 x=-1+2y+\frac{2(1-6y)(1+3y)}{1-2y}Z,
                \qquad (Y,Z)\in[0,1]^2.
\]
This parameterizes the entire displayed range. Multiply each facet
expression by the positive common denominator
\[
 Q(y)=8(1-3y)(1+y)(1-4y)(1+2y)(1-5y)(1+3y)(1-2y)^4.
\]
The resulting polynomial has bidegree at most $(11,4)$.
Every one of these 156 polynomials has been proved nonnegative by
exact Bernstein certificates on dyadic subrectangles.

For a polynomial $P(Y,Z)=\sum a_{ij}Y^iZ^j$ of bidegree $(n,m)$,
the Bernstein coefficient indexed by $(k,l)$ is
\[
 b_{kl}=\sum_{i\le k,j\le l}a_{ij}
          \frac{\binom{k}{i}}{\binom{n}{i}}
          \frac{\binom{l}{j}}{\binom{m}{j}}.
\]
Nonnegative coefficients prove nonnegativity on the unit square.
The script \path{determinant_t4/structural/audit_finite_reduction.py} computes these
coefficients with rational arithmetic and subdivides using the exact
de Casteljau algorithm. Each complete partition and its minimum
coefficient on every box is saved in \path{continuous_hull_N12.json}.
All 156 certificates pass. This proves containment for every
$N\ge12$, $K\ge4$, including spins far beyond any finite search.

\paragraph{The low-$K$ vertices.}
For $K=0$ the map is the identity. For $K=1,2,3$ and $D\ge168$,
the induced map is a convex combination
$\sum_{q=0}^Kp_q\Psi_q$ of local spin-two covariant channels.
The exact coefficients are
\[
 K=1:\quad
 \left(\frac{D-24}{D},\frac{24}{D}\right),
\]
\[
 K=2:\quad
 \left(
 \frac{(D-51)(D-24)}{D(D-3)},\;
 \frac{72(D-24)}{D(D-3)},\;
 \frac{504}{D(D-3)}
 \right),
\]
and
\[
 K=3:\quad\frac1{D(D-8)(D-3)}
 \left(
 \begin{array}{l}
 (D-24)(D^2-131D+3504),\\
 72(D-24)(2D-79),\\
 2520(D-24),\\
 8064
 \end{array}
 \right).
\]
They sum to one and reproduce all four nontrivial multipliers.
Their numerators and denominators have nonnegative coefficients
after substituting $D=168+z$, with strictly positive denominators.
Thus they are nonnegative for every $N\ge12$. These symbolic
identities and positivity statements are checked in
\path{determinant_t4/structural/audit_finite_reduction.py}.

\subsubsection{Small spins require available extensions}
For each $N=5,\ldots,11$, every induced vertex is expressed as an
exact convex combination of local positive maps and listed base
maps whose full spin satisfies $m\le N$.
The decompositions and rational weights are saved in
\path{small_case_certificates.json}.

The extra case $(8,4)$ in the theorem's table is essential for this
particular reduction. It is redundant in the final hull, where a
decomposition can use higher-spin base vectors, but at input spin
$N=8$ one cannot use a spin-nine extension. We therefore retain it
as a separate finite base case. All other needed small-spin vectors
have decompositions respecting $m\le N$.

\subsubsection{From multiplier containment to determinant monotonicity}
Assume the fourteen base-case inequalities hold for every physical
input on their respective full spin sectors. For a full input
$\rho_N$, put $\sigma=R_{N,4}\rho_N$.
If a base map has full spin $m/2$ with $m\le N$, reduce $\rho_N$
to $\rho_m$. This is a physical symmetric spin-$m/2$ state with
the same four-qubit marginal $\sigma$. The assumed base inequality
therefore guarantees that its induced output on $\sigma$ is positive
and has determinant at least $\det\sigma$.
Spin reversal preserves the output spectrum and determinant.

Every local positive, unital, trace-preserving map also cannot
decrease the determinant, because it makes the spectrum more mixed
in the majorization order. Hence every component in the certified
convex decompositions has the required property.
For positive-definite $\sigma$, concavity of log determinant gives
\[
 \log\det\Bigl(\sum_iw_i\tau_i\Bigr)
 \ge\sum_iw_i\log\det\tau_i
 \ge\log\det\sigma.
\]
For singular $\sigma$, the desired inequality follows immediately
from positivity of the target output. The same argument extends
from global channel vertices to their convex mixtures.
The exact decompositions above prove sufficiency for all $N$.
Necessity is immediate, because every listed finite case belongs
to the all-spin claim.

If the unshifted base inequalities are proved, their shifted
versions follow by replacing each full input with
$[\rho+5s\rho_0]/(1+5s)$, whose four-qubit marginal is
$(\sigma+sI_5)/(1+5s)$. Thus the finite reduction also applies
to all positive shifts.

\subsection{Scalar determinant bounds in dimension five}
For a full-rank $d$-dimensional state define
\[
 L_d(\omega)=-\log(d^d\det\omega),\qquad
 V_d(\omega)=\operatorname{Tr}\omega^2-1/d.
\]
Set $L_d=+\infty$ on singular states.

\paragraph{An explicit dimension-five global lower bound.}
For every five-dimensional state,
\begin{equation}\label{t4:eq:d5-burg-lower}
 L_5(\omega)\ge\frac{25}{3}V_5(\omega).
\end{equation}
To prove this minimize $F=L_5-cV_5$, $c=25/3$, over the eigenvalue simplex.
The boundary has $F=+\infty$. Stationarity in the interior gives
$p_i^{-1}+2cp_i=\text{constant}$, with at most two positive roots.
If the larger root occurred twice, varying those coordinates oppositely
would give a negative second variation, contradicting a local minimum.
Thus a nonuniform minimum has spectrum $(a,(1-a)/4,\ldots,(1-a)/4)$,
$a>1/5$. Along this family
\[
 F(a)=-\log\frac{3125a(1-a)^4}{256}
       -\frac{125}{12}(a-1/5)^2,
\qquad
 F'(a)=(5a-1)\left[\frac1{a(1-a)}-\frac{25}{6}\right].
\]
The nonuniform stationary points are $2/5$ and $3/5$. The derivative
is positive, negative, positive on the successive intervals, so only
$a=1/5$ and $a=3/5$ can minimize $F$. The former gives zero, while
\[
 F(3/5)=\log(16/3)-5/3>0.
\]
An exact certificate for the strict inequality is obtained by taking the
first five positive terms in
$\log(16/3)=2\sum_{k\ge0}(13/19)^{2k+1}/(2k+1)$.

\paragraph{A sharp minimum-eigenvalue upper bound.}
For any dimension $d\ge2$, if $\lambda_{\min}(\omega)\ge m$ with
$0<m<1/d$, then
\begin{equation}\label{t4:eq:general-burg-upper}
 L_d(\omega)\le C_d(m)V_d(\omega),\qquad
 C_d(m)=-\frac{d(d-1)}{(1-dm)^2}
 \log\left[\frac{d^d m(1-m)^{d-1}}{(d-1)^{d-1}}\right].
\end{equation}
The constant is sharp at $(m,(1-m)/(d-1),\ldots,(1-m)/(d-1))$.
In particular,
\[
 C_5(m)=-\frac{20}{(1-5m)^2}
          \log\frac{3125m(1-m)^4}{256}.
\]
One proof first establishes the sharp reciprocal-trace bound
\[
 \operatorname{Tr}\eta^{-1}-d^2
 \le\frac{d(d-1)}{u(1-u)}V_d(\eta),
 \qquad\lambda_{\min}(\eta)\ge u.
\]
On each face $p_i\ge u$, stationary free coordinates take at most two
values. The smaller free value cannot occur twice at a maximum.
On the remaining one-small-coordinate family the third derivative is
$-6/b^4+6/[(n-1)^2a^4]<0$, excluding nonuniform interior maxima.
Thus extremizers have $k$ coordinates $u$ and $d-k$ equal coordinates.
Their inverse-excess/variance ratio is
$d(d-k)/[u(1-ku)]$, decreasing in $k$, so $k=1$ is maximal.
Apply this bound to $(\omega+sI)/(1+ds)$ and integrate the resolvent
excess over $s>0$. Equality holds for the same one-small-coordinate
spectrum at every shift, yielding~\eqref{t4:eq:general-burg-upper}.

\paragraph{Sharp lower bound at fixed purity.}
For $0\le V<(d-1)/d$, put
\[
 a_d(V)=\frac1d+\sqrt{\frac{d-1}{d}V},\qquad
 \mathcal L_d(V)=-\log\left[
 \frac{d^d a_d(V)(1-a_d(V))^{d-1}}{(d-1)^{d-1}}\right].
\]
Every state with variance $V$ satisfies $L_d\ge\mathcal L_d(V)$.
To see this, maximize $\sum_i\log p_i$ at fixed trace and purity.
A positive one-large-coordinate spectrum is feasible, so a maximizing
spectrum is interior. At a nonuniform stationary point there are at most
two coordinate values $a>b>0$. The Lagrangian Hessian on a coordinate
of value $a$ has coefficient
$-a^{-2}+(ab)^{-1}>0$. If $a$ occurred twice, the direction that raises
one such coordinate and lowers the other is tangent to both constraints
and has positive second variation. This contradicts a constrained maximum.
Hence exactly one coordinate is large, giving the formula above.
The uniform case follows by continuity; $\mathcal L_d(V)$ is increasing.

\paragraph{Sufficient determinant criterion.}
For physical five-dimensional states $\sigma,\tau$, a variance contraction
$V_5(\tau)\le q^2V_5(\sigma)$ and output floor
$\lambda_{\min}(\tau)\ge m>0$ prove determinant monotonicity whenever
$q^2 C_5(m)\le25/3$. For high-purity inputs the sharper bound
$C_5(m)V_5(\tau)\le\mathcal L_5(V_5(\sigma))$ can be used instead.


\subsection{Sharp spin-two multipole operator norms}
Let $X_\ell$ be Hermitian and supported in tensor rank $\ell$ of a spin-two
operator. Then
\[
 \|X_\ell\|_{\rm op}^2\le c_\ell^2\|X_\ell\|_2^2,
 \qquad
 (c_1^2,c_2^2,c_3^2,c_4^2)
 =\left(\frac25,\frac27,\frac12,\frac{18}{35}\right).
\]
Every constant is sharp.

For rank one, rotate $X_1$ onto $J_z$. Since
$\|J_z\|_{\rm op}^2=4$ and $\operatorname{Tr}J_z^2=10$, the first constant
is $2/5$.

For the even ranks, use a basis in which integer-spin time reversal is
complex conjugation. An even-rank Hermitian operator is real symmetric,
so a real normalized vector attains its largest absolute expectation.
For a pure spin-two state let $n_\ell$ be the squared norm of its tensor-rank
$\ell$ projection, and let $S=|\langle\psi|\Theta\psi\rangle|^2$.
The exact Fierz identity
\[
 n_2=\frac57 n_1+\frac27 S
\]
follows from the invariant two-copy operator
\[
 \sum_m T_{2m}\otimes T_{2m}^\dagger
 -\frac57\sum_m T_{1m}\otimes T_{1m}^\dagger-\frac{10}{7}P_{J=0}.
\]
Its eigenvalues on total spins $J=0,1,2,3,4$ are
$(0,6/7,0,-4/7,0)$, so it vanishes on the symmetric tensor square containing
$|\psi\rangle^{\otimes2}$. Here
$S=5\langle\psi\otimes\psi|P_0|\psi\otimes\psi\rangle$.

For a time-reversal-real pure vector, $n_1=n_3=0$ and $S=1$, hence
$n_2=2/7$. Since $n_0=1/5$ and the pure-state purity is one,
$n_4=18/35$. Therefore, for $\ell=2,4$ and any such real vector,
\[
 |\langle\psi|X_\ell|\psi\rangle|
 \le\|X_\ell\|_2\,\|\mathcal P_\ell(|\psi\rangle\langle\psi|)\|_2
 =c_\ell\|X_\ell\|_2.
\]
This bounds the operator norm. Choosing $X_\ell$ proportional to the
rank-$\ell$ projection of a real pure state attains the bound, proving
sharpness.

For rank three, $X_3$ is purely imaginary antisymmetric in the same real
basis. Its eigenvalues are $\lambda_1,-\lambda_1,\lambda_2,-\lambda_2,0$,
so $\|X_3\|_2^2\ge2\|X_3\|_{\rm op}^2$.
To attain equality, put $u=|2,0\rangle$ and
$v=(|2,2\rangle+|2,-2\rangle)/\sqrt2$, both time-reversal real, and choose
\[
 X_3=\frac{i}{\sqrt2}(|u\rangle\langle v|-|v\rangle\langle u|).
\]
Its rank-one projection vanishes because $\langle u|J_i|v\rangle=0$ for
$i=x,y,z$. It is odd under time reversal, so it is entirely rank three.
Its Hilbert--Schmidt norm is one and its nonzero eigenvalues are
$\pm1/\sqrt2$, establishing the sharp constant $1/2$.



\subsection{Additional physical norm constraints}
Let $\sigma=I_5/5+\sum_{l=1}^4X_l$ be a spin-two density matrix,
with $X_l$ its Hermitian rank-$l$ multipole and $v_l=\Tr X_l^2$.
Set
\[
 a_l=c_l\sqrt{v_l},\qquad
 (c_1^2,c_2^2,c_3^2,c_4^2)=\left(\frac25,\frac27,\frac12,\frac{18}{35}\right).
\]
The constants $c_l$ are the sharp operator-norm constants just proved.
We prove the following necessary conditions:
\begin{align}
8v_2+5v_3+v_4&\le\frac{14}{5},\label{t4:fierzupper}\\
5v_1+4v_4-3v_2&\le\frac65,\label{t4:fierzlower}\\
245a_2^2&\le14+15a_2+20a_4,\label{t4:quadrupole}\\
25a_1^2&\le2+7a_2,\label{t4:variance}\\
15a_1&\le5+2\sqrt5\,a_3.\label{t4:octupole}
\end{align}
Two further total-spin-projector inequalities are
\begin{align}
v_1+v_3&\le\frac15+v_2+v_4,\label{t4:singlet}\\
7v_1+3v_2-8v_3-4v_4&\le\frac{14}{5}.\label{t4:spin2projector}
\end{align}
These conditions hold for every spin-two density matrix, with no
symmetric-extension assumption. Consequently, if a covariant map with
multipliers $g_l$ produces a positive trace-one matrix on the physical
marginal domain, each condition may also be applied after the substitutions
$v_l\mapsto g_l^2v_l$ and $a_l\mapsto |g_l|a_l$.

\subsubsection{The two quadratic Fierz constraints}
Choose Hilbert--Schmidt orthonormal spherical tensors $T_{lq}$ and write
$Q_l=\sum_qT_{lq}\otimes T_{lq}^{\dagger}$. Thus
$\Tr[(\sigma\otimes\sigma)Q_l]=v_l$ for $l>0$.
Let $P_J$ denote the projector onto total spin $J$ in $2\otimes2$.
The exact recoupling identity is
\[
Q_2-\frac57Q_1-\frac{10}{7}P_0
 =\frac67P_1-\frac47P_3.
\]
Put $p_J=\Tr[(\sigma\otimes\sigma)P_J]\ge0$,
$P=\Tr\sigma^2$ and $T=\Tr[\sigma\Theta(\sigma)]$.
The singlet contraction and swap parity give
\[
5p_0=T,\qquad p_1+p_3=\frac{1-P}{2}.
\]
Therefore
\[
-\frac27(1-P)
\le v_2-\frac57v_1-\frac27T
\le\frac37(1-P).
\]
Substitute
$P=1/5+v_1+v_2+v_3+v_4$ and
$T=1/5-v_1+v_2-v_3+v_4$.
The two resulting inequalities are exactly
\eqref{t4:fierzupper}--\eqref{t4:fierzlower}.
The unused projectors provide
\[
p_0=\frac1{25}-\frac15v_1+\frac15v_2-\frac15v_3+\frac15v_4\ge0,
\]
\[
p_2=\frac15-\frac12v_1-\frac3{14}v_2+\frac47v_3+\frac27v_4\ge0,
\]
which are \eqref{t4:singlet}--\eqref{t4:spin2projector}.

\subsubsection{The quadrupole-square constraint}
Write $Q=X_2$ and $u=\Tr Q^2=v_2$.
A real rank-two tensor is equivalent under rotations to a traceless real
symmetric $3\times3$ Cartesian tensor. Diagonalizing that Cartesian tensor
shows that every Hermitian $Q$ can be rotated to
\[
Q=aT_{20}+\frac{b}{\sqrt2}(T_{22}+T_{2,-2}),\qquad a,b\in\mathbb R.
\]
In the $J_z$ basis $(2,1,0,-1,-2)$ this is
\[
Q=\begin{pmatrix}
 2a/\sqrt{14}&0&b/\sqrt7&0&0\\
 0&-a/\sqrt{14}&0&\sqrt{42}b/14&0\\
 b/\sqrt7&0&-2a/\sqrt{14}&0&b/\sqrt7\\
 0&\sqrt{42}b/14&0&-a/\sqrt{14}&0\\
 0&0&b/\sqrt7&0&2a/\sqrt{14}
\end{pmatrix}.
\]
Here $u=a^2+b^2$.
Direct multiplication and orthogonal tensor projection give
\[
P_0(Q^2)=\frac u5I_5,\quad P_1(Q^2)=P_3(Q^2)=0,
\]
\[
\|P_2(Q^2)\|_\HS=\frac{3u}{7\sqrt{14}},\qquad
\|P_4(Q^2)\|_\HS=\frac{12u}{7\sqrt{70}}.
\]
These identities are checked as polynomial identities in $a,b$ by
\path{verify_spin2_norm_constraints.py}.
Since $\Tr\sigma Q=u$, positivity of the variance and Cauchy--Schwarz imply
\begin{align*}
u^2&\le\Tr\sigma Q^2\\
&=\frac u5+\Tr[X_2P_2(Q^2)]+\Tr[X_4P_4(Q^2)]\\
&\le\frac u5+\frac{3u\sqrt u}{7\sqrt{14}}
              +\frac{12u\sqrt{v_4}}{7\sqrt{70}}.
\end{align*}
For $u>0$, divide by $u$. For $u=0$ the desired inequality is automatic.
Substituting $u=7a_2^2/2$ and
$\sqrt{v_4}=\sqrt{35/18}\,a_4$ yields \eqref{t4:quadrupole}.
In particular,
\[
a_4\ge\frac{(245a_2^2-15a_2-14)_+}{20}.
\]

\subsubsection{The dipole constraints}
Rotate the mean spin along the positive $z$ axis and put
$r=\Tr\sigma J_z=|\langle\mathbf J\rangle|$.
Then $v_1=r^2/10$, so $r=5a_1$.
The normalized axial tensors are
\[
T_{20}=\frac1{\sqrt{14}}\operatorname{diag}(2,-1,-2,-1,2),\qquad
T_{30}=\frac1{\sqrt{10}}\operatorname{diag}(1,-2,0,2,-1).
\]
First, $J_z^2=2I+\sqrt{14}T_{20}$, so
\[
r^2\le\Tr\sigma J_z^2\le2+\sqrt{14v_2}=2+7a_2,
\]
which proves \eqref{t4:variance}.
Second,
\[
J_z-\frac{\sqrt{10}}3T_{30}
 =\operatorname{diag}\left(\frac53,\frac53,0,-\frac53,-\frac53\right)
 \le\frac53I.
\]
It follows that
\[
r\le\frac53+\frac{\sqrt{10}}3\sqrt{v_3}
 =\frac53+\frac{2\sqrt5}3a_3,
\]
proving \eqref{t4:octupole}.



\subsection{Dipole lower bound for the determinant deficit}
For every spin-two density matrix $\sigma$, let
$L_5(\sigma)=-\log(5^5\det\sigma)$ and let
$v_1=\Tr[(P_1\sigma)^2]$ denote its squared rank-one multipole norm.
Then
\[
\boxed{\ L_5(\sigma)\ge\frac{25}{2}v_1\ }.
\]
The constant is sharp in the limit of small rank-one perturbations of
$I_5/5$.

Rotate the spin mean along $z$. Write $r=\Tr\sigma J_z\ge0$, so
$r^2=10v_1$. Dephasing in the $J_z$ basis preserves $r$ and increases
the determinant by Hadamard's inequality. It therefore suffices to prove
\[
F(p):=-\sum_{m=-2}^2\log(5p_m)-\frac54\left(\sum_{m=-2}^2mp_m\right)^2
\ge0
\]
for probability vectors $p$. The assertion is immediate if any $p_m=0$.
On the interior, $F$ tends to $+\infty$ at the boundary of the simplex,
since $|r|\le2$. It has a global interior minimum.
At a constrained stationary point,
\[
-\frac1{p_m}-\frac52rm+\lambda=0.
\]
Multiply by $p_m$ and sum to obtain $\lambda=5+(5/2)r^2$. Hence
\[
p_m=\frac{2/5}{2+r^2-rm}.
\]
Every displayed denominator is positive for every real $r$:
$2+r^2-rm\ge2+r^2-2|r|=(|r|-1)^2+1$.
Putting $x=r^2$, a direct rational simplification gives
\[
\sum_{m=-2}^2\frac{2/5}{2+r^2-rm}-1
=-\frac{x^2(5x^3+15x^2+20x+12)}
 {5(x+2)(x^2+4)(x^2+3x+4)}.
\]
For $r\ne0$ the right side is strictly negative, contradicting
normalization. Thus every stationary point has $r=0$, which forces
$p_m=1/5$. At that point $F=0$, proving the bound.

The normalization identity is verified exactly in
\path{verify_dipole_burg_bound.py}.



\subsection{Quadrupole lower bound for the determinant deficit}
\begin{theorem}
For every spin-two density matrix $\sigma$, write
$Q=P_2\sigma$ for its rank-two multipole and $u=\Tr Q^2$.
Then
\[
\boxed{\ -\log(5^5\det\sigma)\ge12u.\ }
\]
\end{theorem}
This bound holds on the entire five-dimensional state space, without a
symmetric-extension assumption.

\paragraph{The spectrum of a Hermitian quadrupole.}
Every $Q$ is rotation-equivalent to
$Q=aT_{20}+b(T_{22}+T_{2,-2})/\sqrt2$, with $a,b$ real and
$u=a^2+b^2$. This follows by diagonalizing the real traceless Cartesian
$3\times3$ tensor corresponding to a rank-two spherical tensor.
Its five eigenvalues consist of
\[
\sqrt{2u/7},\quad-\sqrt{2u/7},
\quad\frac{2a}{\sqrt{14}},
\quad\frac{-a+\sqrt3b}{\sqrt{14}},
\quad\frac{-a-\sqrt3b}{\sqrt{14}}.
\]
If $q$ is the product of the last three, then
\[
q=\frac{a(a^2-3b^2)}{7\sqrt{14}},\qquad
|q|\le\frac{u^{3/2}}{7\sqrt{14}}.
\]
The last inequality follows from
$|a(a^2-3b^2)|\le(a^2+b^2)^{3/2}$, or the triple-angle identity.
Consequently,
\[
\det(\lambda I-tQ)
 =\left(\lambda^2-\frac27t^2u\right)
  \left(\lambda^3-\frac3{14}t^2u\lambda-t^3q\right).
\]
When $t>0$ and $\lambda>t\sqrt{2u/7}$,
\[
\Tr(\lambda I-tQ)^{-1}
=\frac{2\lambda}{\lambda^2-2t^2u/7}
 +\frac{3\lambda^2-3t^2u/14}
 {\lambda^3-3t^2u\lambda/14-t^3q}.
\]
Both denominators and the second numerator are positive. For fixed
$u,\lambda,t$, this expression increases with $q$, so its maximum occurs
at $q=u^{3/2}/(7\sqrt{14})$. This is the axial quadrupole $b=0,a=\sqrt u$.
Writing $z=\sqrt{u/14}$, its eigenvalues are $z(2,-1,-2,-1,2)$.

\paragraph{Exclusion of nonuniform stationary points.}
Consider
\[
F(\sigma)=-\log(5^5\det\sigma)-12\Tr[(P_2\sigma)^2].
\]
It tends to $+\infty$ at the boundary of the density-matrix set, because
the second term is bounded there. Hence it has an interior global minimum.
At any constrained stationary point,
\[
-\sigma^{-1}-24Q+\lambda I=0.
\]
Multiplying by $\sigma$ and taking the trace yields $\lambda=5+24u$,
since $\Tr\sigma Q=u$. Thus
\[
\sigma=(\lambda I-24Q)^{-1},\qquad\lambda=5+24u.
\]
For every $u\ge0$,
\[
\lambda-24\sqrt{2u/7}
=24\left(\sqrt u-\frac12\sqrt{2/7}\right)^2+\frac{23}{7}>0.
\]
The preceding axial comparison is therefore applicable. If $u>0$,
put $z=\sqrt{u/14}>0$; then $\lambda=5+336z^2$ and
\begin{align*}
\Tr\sigma-1
&\le\frac2{336z^2-48z+5}
 +\frac2{336z^2+24z+5}
 +\frac1{336z^2+48z+5}-1\\
&=-\frac{48z^2\left(790272z^4+56448z^3+7392z^2-312z+7\right)}
 {(336z^2-48z+5)(336z^2+24z+5)(336z^2+48z+5)}<0.
\end{align*}
Indeed all denominator factors are positive, the quartic and cubic terms
in parentheses are nonnegative, and
$7392z^2-312z+7>0$ for every real $z$, because
\[
312^2-4\cdot7392\cdot7=-109632<0.
\]
This contradicts $\Tr\sigma=1$. Thus $u=0$ at every stationary point,
which forces $\sigma=I_5/5$. There $F=0$, proving the theorem.

All spectrum, determinant and normalization identities are verified
exactly by \path{verify_quadrupole_burg_bound.py}.



\subsection{Odd-sector lower bound for the determinant deficit}
Let $\Theta$ be integer-spin time reversal on the spin-two state space,
and put $B=(\sigma-\Theta\sigma)/2=X_1+X_3$.
Then every spin-two density matrix satisfies
\[
\boxed{\ -\log(5^5\det\sigma)\ge12\Tr B^2=12(v_1+v_3).\ }
\]
In fact this holds for every five-dimensional conjugation, independently
of the spin interpretation.

Choose a basis in which $\Theta$ is entrywise conjugation. Then $B$ is
purely imaginary Hermitian antisymmetric, with eigenvalues in opposite
pairs and at least one zero eigenvalue.
Consider
\[
F(\sigma)=-\log(5^5\det\sigma)-12\Tr B^2.
\]
This tends to $+\infty$ at the boundary of the state space, so its minimum
is attained in the interior. The derivative of $\Tr B^2$ with respect to
$\sigma$ is $2B$, since the odd projection is self-adjoint for the
Hilbert--Schmidt inner product. Thus at any constrained stationary point,
\[
\sigma^{-1}=\lambda I-24B,\qquad
\lambda=5+24\Tr B^2\ge5.
\]
Time reversal gives $(\Theta\sigma)^{-1}=\lambda I+24B$.
The identity $B=(\sigma-\Theta\sigma)/2$ therefore implies
\[
B=24B(\lambda^2I-576B^2)^{-1}.
\]
Every nonzero eigenvalue $b$ of $B$ obeys
$\lambda^2-576b^2=24$. Thus all nonzero eigenvalues have the same magnitude,
and there are $k=1$ or $2$ nonzero opposite pairs. Write their squared
magnitude as $b^2$. Combining
$\Tr B^2=2kb^2$, $\lambda=5+48kb^2$ and the preceding equation yields
\[
k\lambda^2-12\lambda+12(5-2k)=0.
\]
For $k=2$, this gives $\lambda=3\pm\sqrt3<5$, impossible.
For $k=1$, it gives $\lambda=6$, $b^2=1/48$ and $\Tr B^2=1/24$.
The stationary matrix has two eigenvalues $(6\pm\sqrt{12})^{-1}$ and
three eigenvalues $1/6$. Hence
\[
\det\sigma=\frac1{5184},\qquad
F(\sigma)=\log\frac{5184}{3125}-\frac12>0.
\]
For an exact verification of the last strict inequality, let
$r=2059/8309$. Then
\[
\log\frac{5184}{3125}=2\operatorname{arctanh}r
 >2\left(r+\frac{r^3}{3}\right)>\frac12,
\]
where the second inequality is rational and is checked in
\path{verify_odd_burg_bound.py}.
The remaining stationary possibility is $B=0$, which forces
$\sigma=I_5/5$ and $F=0$. These cases exhaust all stationary points, so
the global minimum is zero, as claimed.



\subsection{Nonlinear lower bounds near sector boundaries}
Let $\sigma$ be a spin-two density matrix, let
$L_5(\sigma)=-\log(5^5\det\sigma)$, and write $v_l=\Tr[(P_l\sigma)^2]$.
The following lower bounds hold on the entire five-dimensional state space:
\begin{align*}
L_5(\sigma)&\ge B_Q\left(\sqrt{7v_2/2}\right),\\
B_Q(p)&=-3\log(1-p)-4\log(1+p)+2\log(1+p/2),\\[2mm]
L_5(\sigma)&\ge B_D\left(\sqrt{5v_1/2}\right),\\
B_D(p)&=\log(1-p^2/4)-4\log(1-p^2),\\[2mm]
L_5(\sigma)&\ge-4\log\bigl(1-2(v_1+v_3)\bigr).
\end{align*}
The arguments $p$ lie in $[0,1]$, and $v_1+v_3\le1/2$.
Each expression is zero at zero and is interpreted as $+\infty$ at
its upper endpoint. The bounds can be combined by taking their maximum,
also with the previously proved linear lower bounds.

\subsubsection{A general determinant bound}
Let $H$ be Hermitian with $\|H\|_{\mathrm{op}}\le1$, let
$h=\Tr\sigma H$, and let $0\le x<1$. The matrix $A=I-xH$ is positive
definite. Matrix arithmetic--geometric mean, applied to
$A^{1/2}\sigma A^{1/2}$, gives
\[
\det A\det\sigma
\le\left(\frac{\Tr A\sigma}{5}\right)^5
=\left(\frac{1-xh}{5}\right)^5.
\]
Therefore
\[
\boxed{\quad L_5(\sigma)\ge\log\det(I-xH)-5\log(1-xh).\quad}
\]
No tracelessness assumption on $H$ is needed for this general statement.
The formula also follows by optimizing the scalar multiplier in the
standard variational formula for negative log determinant.

\subsubsection{The quadrupole bound}
Put $Q=P_2\sigma$, $u=\Tr Q^2$, and suppose $u>0$.
Set $b=\sqrt{2u/7}$ and $H=Q/b$. The exact quadrupole spectrum established
in the preceding quadrupole argument gives $\|H\|_{\mathrm{op}}=1$.
It consists of $1,-1$ and three real numbers of sum zero, squared sum
$3/2$, and product $q$ with $|q|\le1/4$. Consequently,
\begin{align*}
\det(I-xH)&=(1-x^2)(1-3x^2/4-qx^3)\\
&\ge(1-x^2)(1-3x^2/4-x^3/4)\\
&=(1-x)^2(1+x)(1+x/2)^2.
\end{align*}
Also
\[
h=\frac{\Tr\sigma Q}{b}=\frac ub=\sqrt{7u/2}=:p\in[0,1],
\]
because $\Tr\sigma Q=u$ and $h\le\|H\|_{\mathrm{op}}=1$.
For $0<p<1$, choose $x=p$ in the general determinant bound. This gives
\[
L_5(\sigma)\ge
2\log(1-p)+\log(1+p)+2\log(1+p/2)-5\log(1-p^2)
=B_Q(p).
\]
For $u=0$, the bound is $L_5\ge0$.
For $p=1$, retain $x<1$ and let $x\uparrow1$; the bound tends to
$+\infty$, proving that the state is singular.

\subsubsection{The dipole bound}
Rotate the mean spin along $z$ and write $r=\Tr\sigma J_z\ge0$.
Set $H=J_z/2$ and $p=r/2=\sqrt{5v_1/2}\in[0,1]$.
Its eigenvalues are $1,1/2,0,-1/2,-1$, so
\[
\det(I-xH)=(1-x^2)(1-x^2/4).
\]
Taking $x=p$ gives
\[
L_5(\sigma)\ge\log(1-p^2/4)-4\log(1-p^2)=B_D(p).
\]
The endpoints follow by the same zero-value and limiting arguments.

\subsubsection{The time-reversal-odd bound}
Write $B=(\sigma-\Theta\sigma)/2=P_1\sigma+P_3\sigma$ and
$w=\Tr B^2=v_1+v_3$. For $w>0$, put $H=B/\sqrt{w/2}$.
The eigenvalues of $B$ are $b_1,-b_1,b_2,-b_2,0$, where
$b_1^2+b_2^2=w/2$. Thus $\|H\|_{\mathrm{op}}\le1$ and its nonzero
squared eigenvalue magnitudes are $a,1-a$ for some $a\in[0,1]$.
It follows that
\[
\det(I-xH)=(1-ax^2)(1-(1-a)x^2)
=1-x^2+a(1-a)x^4\ge1-x^2.
\]
Since $\Tr\sigma B=w$, we have
\[
h=\frac w{\sqrt{w/2}}=\sqrt{2w}=:p\in[0,1].
\]
Choosing $x=p$ gives
\[
L_5(\sigma)\ge-4\log(1-p^2)=-4\log(1-2w).
\]
Again $w=0$ is immediate, while $w=1/2$ follows by a limiting argument
and forces singularity.

All algebraic determinant identities are verified exactly in
\path{verify_nonlinear_burg_bounds.py}.

\subsubsection{Optional optimization of the scalar parameter}
The same proof permits any $0\le x<1$. For the quadrupole, the stronger
one-dimensional lower bound is
\[
L_5(\sigma)\ge\max_{0\le x<1}
\{2\log(1-x)+\log(1+x)+2\log(1+x/2)-5\log(1-px)\}.
\]
For $0<p<1$, its unique maximizer is
\[
x_Q(p)=\frac{5p-7+\sqrt{49+130p+145p^2}}{10+6p}.
\]
Indeed the sign of the derivative is the negative of
$(5+3p)x^2+(7-5p)x-10p$, whose derivative is positive on $[0,1]$
and whose endpoint values have opposite signs. The polynomial is negative
at $x=p$, so $p<x_Q(p)<1$.
For the odd sector the analogous bound is
\[
L_5(\sigma)\ge\max_{0\le x<1}
\{\log(1-x^2)-5\log(1-px)\}.
\]
Its unique maximizer is
\[
x_O(p)=\frac{5p}{1+\sqrt{1+15p^2}},
\]
since the derivative changes sign at the unique positive root of
$3px^2+2x-5p$. Again $p<x_O(p)<1$.
The zero endpoint is obtained by continuity. The simpler explicit bounds
above correspond to the admissible choice $x=p$.


\subsection{Exact box certificates for eleven order-four base channels}
This section documents the finite certificates in
\path{bounds/norm_floor_certificate_N*_K*.json}.
They establish determinant monotonicity for eleven of the fourteen base
channels in the exact all-spin finite reduction. No numerical optimization
result is used as a proof: optimization only proposes rational identities,
which are then checked exactly, and every final scalar comparison is rational.

\paragraph{Coordinates and physical positive maps.}
Write a physical four-body marginal as
\[
 \sigma=I_5/5+\sum_{l=1}^4X_l,\qquad
 v_l=\operatorname{Tr}X_l^2,\qquad a_l=c_l\sqrt{v_l},\qquad
 (c_1^2,c_2^2,c_3^2,c_4^2)=(2/5,2/7,1/2,18/35).
\]
The operator-norm estimates $\|X_l\|_{\rm op}\le a_l$ were proved above.
The target output is $\tau=I_5/5+\sum_l f_lX_l$.
Let $g^{(j)}$ be a finite list of multiplier vectors such that
$\eta_j=I_5/5+\sum_lg_l^{(j)}X_l$ is a density matrix for every physical
input marginal. These maps need only be positive on the physical marginal
domain; no determinant monotonicity assumption about them is required.
The list used here consists of the five irreducible covariant channels
from spin $N/2$ to spin two, their time reversals, and their compositions
with the three local positive maps below and their time reversals.
For odd $N$, the global even-dimensional Breuer--Hall map also gives the
columns $g_l^{(j)}h_l$, where
$h=(0,-2/(N-1),0,-2/(N-1))$.

For completeness, all three additional local maps arise by twirling
\[
 X\longmapsto\operatorname{Tr}(X)I-X-VX^TV^\dagger,
 \qquad V=A U_{\rm TR},\quad V^T=-V,\quad \|V\|_{\rm op}\le1.
\]
For rank-one inputs, the vectors $\psi$ and $V\overline\psi$ are
orthogonal and the latter has norm at most one, proving positivity.
After twirling, normalize by $4-\nu$, where
$\nu=\operatorname{Tr}AA^\dagger/5$. The normalized maps are unital and
trace preserving. The three choices are
\[
 A_1=J_+/\sqrt6,\qquad A_2=\operatorname{diag}(1,1,0,-1,-1),
 \qquad
 A_3=\begin{pmatrix}
 0&1&0&0&0\\0&0&2\sqrt6/5&0&0\\0&0&0&2\sqrt6/5&0\\
 0&0&0&0&1\\0&0&0&0&0
 \end{pmatrix}.
\]
Writing $\nu_1,\nu_3$ for their rank-one and rank-three squared
Hilbert--Schmidt norms divided by five, the pairs are respectively
\[
 (2/3,0),\qquad(18/25,2/25),\qquad(484/625,6/625).
\]
Their rank-$l$ multipliers are
\[
 b_l=\frac{-1-(-1)^l\{\nu_1 f_l(4,1)+\nu_3 f_l(4,3)\}}
 {4-\nu_1-\nu_3}.
\]
To justify the multiplier formula, consider the twirl of $X\mapsto AXA^\dagger$.
The vectorized $A$ has rank-one and rank-three components only. Averaging
its Choi projector over rotations removes cross terms and replaces each
rank-$L$ diagonal block by its trace divided by $2L+1$. Comparison with
the vertex Choi normalization in Section~2.2 gives exactly
$\nu_1\Phi_1^{(4)}+\nu_3\Phi_3^{(4)}$. The other terms of the twirled
positive map are $\operatorname{Tr}(X)I-X$ and the negative of that
map composed with $\Theta$, proving the displayed $b_l$. Its value on
$I$ and its trace multiplier are both $4-\nu_1-\nu_3>0$, which proves
the stated normalization. The three matrices displayed above satisfy
$V^{\mathsf T}=-V$ and $VV^\dagger\preceq I$ by direct multiplication;
projection onto $J_z,J_+$ and their adjoints gives the listed $\nu$ pairs.

The exact Racah multipliers and global channel multipliers are reconstructed
by the verifier, independently of their values stored in each certificate.

\paragraph{A certified lower bound for the output spectrum.}
Suppose nonnegative rational coefficients satisfy
\[
 f_l=\sum_jw_jg_l^{(j)}+h_l^+-h_l^-,\qquad l=1,\ldots,4.
\]
Then
\[
 \tau=\sum_jw_j\eta_j+
 \frac{1-\sum_jw_j}{5}I+
 \sum_l(h_l^+-h_l^-)X_l
\]
and therefore
\[
 \lambda_{\min}(\tau)\ge
 m(a):=\frac{1-\sum_jw_j}{5}-\sum_l(h_l^++h_l^-)a_l.
\]
This affine bound decreases in each coordinate. On a box
$a_l\in[\underline a_l,\overline a_l]$, the single rational value
$m(\overline a)$ is thus a valid uniform spectral floor.
The stored floor bases contain every coefficient of this exact identity.

\paragraph{Initial domain and exclusions.}
For each known positive image, $\operatorname{Tr}\eta_jX_l'=\|X_l'\|_2^2$
and the operator-norm estimate imply
$(g_l^{(j)})^2v_l\le c_l^2$. Hence the initial domain is the rational box
\[
 0\le a_l\le c_l^2/\max_j|g_l^{(j)}|,
\]
where the identity image is included. Each image also obeys
$\sum_l(g_l^{(j)})^2v_l\le4/5$. The two further linear inequalities
\[
 8v_2+5v_3+v_4\le14/5,\qquad 5v_1+4v_4-3v_2\le6/5,
\]
and the three inequalities
\[
 245a_2^2-15a_2-14\le20a_4,\quad
 25a_1^2\le2+7a_2,\quad15a_1\le5+2\sqrt5a_3
\]
hold for every spin-two state; as proved in the physical norm constraints above.
They can be applied to every positive image by replacing $v_l$ by
$(g_l^{(j)})^2v_l$ and $a_l$ by $|g_l^{(j)}|a_l$.
A box is excluded only when the exact minimum of one left-hand side on
that box exceeds its right-hand side. The last square-root inequality
is squared only when its left-hand side is positive.
If the input is forced to be singular, determinant monotonicity holds
immediately because the output is positive semidefinite.

\paragraph{Scalar comparisons on the remaining boxes.}
Put $L(\omega)=-\log(5^5\det\omega)$ and $V=\sum_lv_l$.
For a uniform floor $m>0$,
\[
 L(\tau)\le C_5(m)\sum_lf_l^2v_l,\qquad
 C_5(m)=\frac{20}{(1-5m)^2}
 \log\frac{256}{3125m(1-m)^4}.
\]
The global and sector lower bounds used for the input are
\[
 \begin{split}
 L(\sigma)\ge\max\bigl\{&(25/3)V,(25/2)v_1,12v_2,12(v_1+v_3),\\
 &\mathcal L_5(V),B_D((5/2)a_1),B_Q((7/2)a_2),\\
 &-4\log(1-2(v_1+v_3))\bigr\},
 \end{split}
\]
where
\[
 \mathcal L_5(V)=\log\frac{256}{3125A(1-A)^4},\quad
 A=1/5+\sqrt{4V/5},
\]
\[
 B_D(p)=\log\frac{1-p^2/4}{(1-p^2)^4},\qquad
 B_Q(p)=\log\frac{(1+p/2)^2}{(1-p)^3(1+p)^4}.
\]
The dimension-five scalar, dipole, quadrupole, odd-sector and nonlinear
proofs have all been given above. The nonlinear expressions are increasing on their
domains, so evaluation at lower box endpoints is valid.

A box certificate may use a ratio bound, a direct comparison of upper
and lower endpoint values, or a convex combination of the four linear
input lower bounds. For the latter, rational weights $\lambda_i\ge0$,
$\sum_i\lambda_i=1$, produce a linear lower bound with coefficients
$d_l=\sum_i\lambda_i d_l^{(i)}$. It suffices to verify
\[
 \sum_l\min\{(d_l-C_5(m)f_l^2)\underline v_l,
                  (d_l-C_5(m)f_l^2)\overline v_l\}\ge0.
\]
This is an exact sufficient condition on the entire box.
Linear-fractional ratios are maximized at box vertices; when the
denominator is the maximum of $(25/3)V$ and $(25/2)v_1$, intersections
of their equality plane with box edges are included as well.

\paragraph{Exact arithmetic and coverage.}
Logarithms are enclosed by rational partial sums of
$\log x=2\sum_{k\ge0}z^{2k+1}/(2k+1)$, $z=(x-1)/(x+1)$,
after power-of-two range reduction; the omitted positive tail is bounded
by its geometric majorant. Twelve terms suffice for the saved comparisons.
Square roots in $\mathcal L_5$ are bounded below using an integer square
root on the grid $2^{-40}\mathbb Z$.
The root box is bisected along successive coordinates. The verifier
reconstructs its full binary subdivision tree from the rational endpoints,
checks that every internal node has both children and that leaves do not
overlap, and then verifies every exact predicate. Thus coverage does not
rely on an optimization status flag or a numerical volume sum.

\begin{center}
\begin{tabular}{rrr}
$N$&$K$&Certified leaf boxes\\\hline
5&4&353\\
6&4&219\\
6&5&16\\
7&4&48\\
7&5&36\\
7&6&15\\
8&4&41290\\
8&7&16\\
9&8&46\\
10&9&397\\
11&10&12567
\end{tabular}
\end{center}
For each of these eleven cases the saved partition is complete and the
independent verifier succeeds. The remaining finite cases $(5,2)$,
$(5,3)$ and $(6,3)$ are not proved by these certificates.

\subsection{Sharp output Burg bound from a floor and a purity bound}
For $0<m<1/5$ and $s\ge0$, consider the compact eigenvalue set
\[
 \mathcal E(m,s)=\{p_i\ge m:\ \sum_{i=1}^5p_i=1,
                         \ \sum_{i=1}^5p_i^2\le1/5+s\}.
\]
The following finite formula gives the exact maximum of
$L(p)=-\sum_i\log(5p_i)$ on this set.
If $s\ge(4/5)(1-5m)^2$, the maximum is attained at
$(1-4m,m,m,m,m)$, and equals
\[
 U(m,s)=\log\frac1{3125m^4(1-4m)}.
\]
Otherwise, for each $k=0,1,2,3$, set
\[
 q=5-k,\quad S=1-km,\quad Q=1/5+s-km^2,\quad
 \Delta=qQ-S^2,
\]
\[
 b=\frac{S-\sqrt{(q-1)\Delta}}q,\qquad
 a=\frac{S+\sqrt{\Delta/(q-1)}}q.
\]
Retain the candidate when $\Delta\ge0$ and $b\ge m$.
Then
\[
 \boxed{\quad U(m,s)=
 \max_{\text{retained }k}\log\frac1{3125m^kba^{q-1}}.\quad}
\]
The case $s=0$ has $U=0$. At least one candidate is retained whenever
the vertex branch is not applicable.

\paragraph{Proof.}
The simplex vertex has maximal purity $(1-4m)^2+4m^2$ and is a
majorization maximum on the floor-constrained simplex, proving the first
branch. Suppose it is excluded by the purity constraint.
At a maximizer, take the smallest face of the floor-constrained simplex
containing that maximizer, so precisely $k$ entries equal $m$ and its
remaining $q$ entries are strictly larger than $m$.
If the purity constraint were inactive and $q\ge2$, a nonzero direction
with zero coordinate sum would be feasible in both signs locally.
The second derivative of $-\sum_i\log p_i$ in that direction is strictly
positive, contradicting a local maximum. The alternative $q=1$ is the
excluded simplex vertex. Thus purity is active.

If all free coordinates are equal, then $\Delta=0$ and the displayed
formula already includes the candidate. Otherwise trace and purity
constraints have independent gradients on the face. Stationarity gives
\[
 -p_i^{-1}=\lambda+2\mu p_i.
\]
There are at most two positive free-coordinate values, say $a>b$.
Subtracting their equations gives $2\mu=1/(ab)$. The Lagrangian Hessian
coefficient at $b$ is
\[
 b^{-2}-(ab)^{-1}>0.
\]
If $b$ occurred at least twice, the direction increasing one such entry
and decreasing the other would be tangent to both constraints, and have
positive second variation. This contradicts the second-order condition
at a constrained maximum. Hence the smaller value occurs once and the
larger value $q-1$ times. The trace and purity equations give the stated
$a,b$. If a candidate has $b=m$, it belongs to the boundary of the chosen
face and is harmlessly counted twice. Every global maximum has now been
included, and each retained candidate is feasible, proving equality.

\paragraph{Rational interval evaluation.}
The certificate verifier never needs to evaluate the displayed roots or
logarithms in floating point. Rational lower and upper bounds on each
root give $b_-$, $b_+$ and $a_-$. A candidate is skipped only if
$\Delta<0$ or $b_+<m$. Otherwise, even when feasibility is uncertain,
\[
 p_1p_2p_3p_4p_5\ge m^k\max(m,b_-)\max(m,a_-)^{q-1}
\]
for that candidate. Applying an upper rational logarithm enclosure gives
a conservative upper bound. Taking the maximum over all unexcluded
candidates is therefore rigorous.

\subsection{Two compact-domain K=3 certificates}
The files
\path{compact_norm_floor_certificate_N5_K3.json} and
\path{compact_norm_floor_certificate_N6_K3.json} establish
$L(\tau)\le L(\sigma)$ on the respective domains
\[
 N=5:\quad p_2=\sqrt{7v_2/2}\le31/32,
 \qquad
 N=6:\quad p_2\le511/512.
\]
The first has 488 leaf boxes and five exact output-floor identities;
the second has 696 leaf boxes and six identities. Both complete
partitions passed the independent exact verifier.
Their complementary endpoint regions are handled separately by the
analytic endpoint estimates, so these compact certificates are not
presented as proofs on the entire initial box.

On a box, the spectral floor is evaluated at its upper corner and
$s=\sum_lf_l^2\overline v_l$. The above $U(m,s)$ is compared to lower
bounds on the input at the box's lower corner. Besides the earlier
bounds, an arbitrary stored rational parameter $x\in[0,1)$ gives
\[
 L(\sigma)\ge
 \log\frac{(1-x)^2(1+x)(1+x/2)^2}{(1-xp_2)^5}.
\]
For the odd sector $w=v_1+v_3$, a stored rational $t\ge0$ satisfying
$2wt^2<1$ and $2wt<1$ gives
\[
 L(\sigma)\ge\log\frac{1-2wt^2}{(1-2wt)^5}.
\]
These follow from the general determinant arithmetic--geometric-mean
bound proved in the nonlinear bounds above; the second uses
$x=t\sqrt{2w}$. At the lower endpoint $w_0$ of a box, the selected
$x=t\sqrt{2w_0}$ is fixed when invoking monotonicity in the expectation
parameter, so the endpoint evaluation is a valid bound throughout the
box. No optimality claim about the stored parameters is required.
Both parameter choices are rationally checked, and the logarithm of the
resulting rational expression is enclosed exactly.


\subsection{Completing both K=3 cases at the boundary}
Let $\rho$ be a spin-$N/2$ state for $N=5$ or $6$, and put
$\sigma=R_{N,4}\rho$, $\tau=R_{N,4}\Phi_3^{(N)}\rho$,
$u=\Tr[(P_2\sigma)^2]$ and $p=\sqrt{7u/2}$.
Then $0\le p\le1$.
\begin{theorem}
The determinant inequality $\det\tau\ge\det\sigma$ holds throughout
\[
N=5:\quad p\ge\frac{31}{32},\qquad
N=6:\quad p\ge\frac{511}{512}.
\]
For $p<1$ these certificates give the strict inequality
$L_5(\tau)<L_5(\sigma)$, where $L_5=-\log(5^5\det)$.
\end{theorem}

\subsubsection{Exact positive-output decompositions}
Let $G_q$ denote the extremal covariant channel from spin $N/2$ to spin
$2$ with environment spin $(N-4)/2+q$; $G_0=R_{N,4}$.
Each $G_q(\rho)$ is positive with trace one.
Writing $X_2=P_2\sigma$, exact multiplier comparison gives
\begin{align*}
N=5:\quad\tau&=\frac4{63}I_5+\frac4{105}G_0(\rho)
 +\frac{29}{45}G_2(\rho)-\frac29X_2,\\
N=6:\quad\tau&=\frac1{45}I_5+\frac1{75}G_0(\rho)
 +\frac{19}{45}G_1(\rho)+\frac{34}{75}G_2(\rho)-\frac7{90}X_2.
\end{align*}
The constant terms and channel weights preserve trace in both identities.
The sharp quadrupole norm identity is
$\|X_2\|_{\mathrm{op}}=\sqrt{2u/7}=2p/7$.
Consequently, with $\delta=1-p$,
\[
N=5:\quad\lambda_{\min}(\tau)\ge\frac4{63}\delta,
\qquad
N=6:\quad\lambda_{\min}(\tau)\ge\frac1{45}\delta.
\]

\subsubsection{Uniform output-purity bounds}
The exact bounds needed below are
\[
N=5:\quad\Tr\tau^2\le\frac{6739}{22050},\qquad
N=6:\quad\Tr\tau^2\le\frac13.
\]
Here is an explicit construction of the two-copy operator used for this
table. Let $\Lambda=R_{N,4}\Phi_3^{(N)}$, let $\Lambda^*$ be its trace
adjoint, and let $E_{ab}=|a\rangle\langle b|$ in the five-dimensional
output basis. Put
\[
 O_N=\sum_{a,b=0}^4\Lambda^*(E_{ab})\otimes\Lambda^*(E_{ba}).
\]
Expanding the trace gives
$\operatorname{Tr}[(\rho\otimes\rho)O_N]=\operatorname{Tr}\Lambda(\rho)^2$.
Covariance makes $O_N$ invariant under simultaneous rotations, so it is
scalar on each total-spin copy. For its projector
\[
 P_J=\sum_{M=-J}^J|J,M\rangle\langle J,M|,\qquad
 |J,M\rangle=\sum_{m,n}\langle N/2,m;N/2,n\mid J,M\rangle|m,n\rangle,
\]
the scalar is $\operatorname{Tr}(P_JO_N)/(2J+1)$. Use the reduction
Kraus matrices (5), the square channel Kraus construction in Section~2.4,
and the Clebsch--Gordan recurrence there in this finite formula. Its eigenvalues
on the symmetric tensor square of the global spin space are
\[
\begin{array}{c|c|c}
N&J&\text{purity-operator eigenvalue}\\\hline
5&1&11413/61250\\
 &3&62291/275625\\
 &5&6739/22050\\\hline
6&0&-1/18\\
 &2&19/150\\
 &4&16/81\\
 &6&1/3
\end{array}
\]
A pure tensor $|\psi\rangle^{\otimes2}$ belongs to the symmetric tensor
square, so its expectation is bounded by the maximum listed eigenvalue.
Convexity of purity extends this upper bound to mixed input states.
Equivalently, the embedded checker evaluates the same contractions
using exact recoupling coefficients in
\path{verify_k3_endpoint.py}.
Thus $V_{\rm out}=\Tr\tau^2-1/5$ obeys
\[
N=5:\quad V_{\rm out}\le\frac{2329}{22050},\qquad
N=6:\quad V_{\rm out}\le\frac2{15}.
\]

\subsubsection{A common input lower bound}
The nonlinear quadrupole bound gives, for $0<\delta<1$,
\[
L_5(\sigma)\ge3\log(1/\delta)-4\log(2-\delta)
 +2\log((3-\delta)/2).
\]
For $\delta\le1/32$, use
$\log(2-\delta)\le\log2\le7/10$ and
\[
\log((3-\delta)/2)\ge\log(95/64)
\ge2\frac{95-64}{95+64}=\frac{62}{159}.
\]
Writing $\ell=\log(1/\delta)$, this yields
\[
L_5(\sigma)\ge3\ell-\frac{1606}{795}.
\]
Here $\log2\ge2/3$ follows from the same elementary bound
$\log x\ge2(x-1)/(x+1)$ for $x\ge1$.
The upper bound $\log2<7/10$ follows from
$\sum_{k=0}^3(7/10)^k/k!>2$.

\subsubsection{The $N=5$ endpoint}
Let $\delta\le1/32$ and $m=(4/63)\delta\le1/504$.
The established minimum-eigenvalue Burg upper bound gives
\[
L_5(\tau)\le\frac{20V_{\rm out}}{(1-5m)^2}
 \log\frac{256}{3125m(1-m)^4}.
\]
Since $20V_{\rm out}\le4658/2205$, $\log x\le x-1$ and
$-\log(1-m)\le m/(1-m)$ give
\[
L_5(\tau)\le A_5(\ell+b_5),\qquad
A_5=\frac{4658}{2205}\left(\frac{504}{499}\right)^2
=\frac{2683008}{1245005}<3,
\]
\[
b_5=\frac{907}{3125}+\frac4{503}.
\]
Indeed the constant in the logarithm is
$256/[3125(4/63)]=4032/3125$.
Also $\ell\ge5\log2\ge10/3$. Therefore
\begin{align*}
L_5(\sigma)-L_5(\tau)
&\ge(3-A_5)\ell-A_5b_5-\frac{1606}{795}\\
&\ge\frac{47878982324888}{311161765265625}>0.
\end{align*}

\subsubsection{A determinant bound using the output purity}
For a five-dimensional state with $\lambda_{\min}\ge m$,
$\Tr\tau^2\le1/3$, and $0<m\le1/400$, we claim
\[
\det\tau\ge\frac{m^2}{64}.
\]
Here is a direct derivation of the needed extremal spectrum.
Minimize the logarithm of the product of the five eigenvalues subject to
$p_i\ge m$, $\sum p_i=1$ and $\sum p_i^2\le1/3$.
If the purity constraint were inactive, strict concavity along any
trace-preserving variation of two free coordinates would exclude a
minimum with two or more coordinates greater than $m$.
The only remaining candidate has four coordinates equal to $m$, but its
purity $(1-4m)^2+4m^2$ exceeds $1/3$ in the range in question.
Thus the purity constraint is active.
On a face with $k$ coordinates equal to $m$, an all-equal free spectrum
would have purity $km^2+(1-km)^2/(5-k)$, which is not $1/3$ in this
range. The trace and purity gradients on the free coordinates are therefore
independent. At an interior constrained extremum the free coordinates
have at most two values $a>b$:
the multiplier equation is $1/p_i=\lambda+2\nu p_i$.
The smaller value cannot be repeated. Indeed $2\nu=-1/(ab)$, and the
constrained second derivative on a variation of two equal smaller
coordinates is proportional to $-1/b^2+1/(ab)<0$.
Consequently, the free spectrum has one $b$ and $4-k$ copies of $a$.
The cases $k=0,1$ give a negative $b$. For $k=0$ this is immediate;
for $k=1$,
\[
b=\frac{1-m-\sqrt{1+6m-15m^2}}4<0,
\]
because $1+6m-15m^2-(1-m)^2=8m(1-2m)>0$.
For $k=3$, the residual squared-displacement condition is
$-1/3+6m-15m^2\ge0$, which fails for $m\le1/400$.
The case $k=4$ is infeasible as already observed.
Thus the minimum product has spectrum $(m,m,b,a,a)$, where
\[
b=\frac{1-2m-\sqrt{8m-20m^2}}3,\qquad
 a=\frac{1-2m+\sqrt{2m-5m^2}}3.
\]
In this range $a\ge(1-2m)/3>1/4$.
Also $b\ge1/4$, since $1/4-2m>0$ and
\[
(1/4-2m)^2-(8m-20m^2)
=\frac1{16}-9m+24m^2\ge\frac1{16}-\frac9{400}>0.
\]
This proves $\det\tau=m^2ba^2\ge m^2/64$ at the minimum, hence everywhere.

\subsubsection{The $N=6$ endpoint}
Let $\delta\le1/512$ and $m=\delta/45<1/400$.
The preceding determinant bound and $\Tr\tau^2\le1/3$ give
\[
L_5(\tau)\le-\log\frac{3125m^2}{64}
=2\ell+\log\frac{5184}{125}<2\ell+\frac{15}{4}.
\]
The last logarithm bound is exact: the Taylor polynomial satisfies
\[
\sum_{k=0}^8\frac{(15/4)^k}{k!}-\frac{5184}{125}
=\frac{3077585521}{7340032000}>0.
\]
Also $\ell\ge9\log2>6$. Thus
\[
L_5(\sigma)-L_5(\tau)
>\ell-\frac{15}{4}-\frac{1606}{795}
>6-\frac{15}{4}-\frac{1606}{795}
=\frac{731}{3180}>0.
\]
At $p=1$, the input nonlinear Burg bound diverges, so $\det\sigma=0$
and the determinant comparison follows immediately from positivity of
$\tau$. This completes both endpoint certificates.




\subsection{The exceptional base case (N,K)=(5,2)}
\begin{theorem}
Let $\Phi_2$ be the extremal SU(2)-covariant channel from spin $5/2$ to itself
with environment spin $K=2$, and let $R=R_{5,4}$ denote reduction from the
symmetric five-qubit space to the symmetric four-qubit space. Then
\[
 \boxed{\det R\Phi_2(\rho)\ \geq\ \det R(\rho)\qquad(\rho\succeq0).}
\]
In particular, the unshifted determinant inequality holds for every mixed
spin-$5/2$ input, with arbitrary simultaneous complex coherences.
\end{theorem}

This is a computer-assisted exact weighted sum-of-squares proof. The embedded
rational data specify the polynomial matrices and Gram matrices. The verifier
checks the channel, covariance, polynomial identity, and positivity using
integer and rational arithmetic. An SDP solver is not needed to verify the
certificate.

\subsubsection{Channel and rational coordinates}
Use the Dicke basis indexed by $a=0,\ldots,5$, and the Kraus convention
\[
 (K_q)_{ra}=\left\langle \tfrac52,\tfrac52-a;2,q
       \mathrel{\big|}\tfrac52,\tfrac52-r\right\rangle,
 \qquad q=-2,\ldots,2.
\]
Thus $\Phi_2(\rho)=\sum_qK_q\rho K_q^\dagger$ and
$\sum_qK_q^\dagger K_q=I_6$. The two reduction Kraus operators satisfy
\[
 (F_b)_{r,r+b}=\sqrt{\frac{\binom4r}{\binom5{r+b}}},\qquad b=0,1.
\]
Put $D_n=\diag\bigl(\sqrt{\binom n0},\ldots,\sqrt{\binom nn}\bigr)$ and
$\rho=D_5HD_5$. Then $\rho\succeq0$ if and only if $H\succeq0$.
Writing $R(\rho)=D_4A(H)D_4$ and $R\Phi_2(\rho)=D_4B(H)D_4$ gives
\[
 A(H)_{rs}=H_{rs}+H_{r+1,s+1},\qquad
 \Delta(H)=96\bigl(\det B(H)-\det A(H)\bigr).
\]
The linear map $B$ has rational coefficients. The file
\path{coherent_sos/channel_rational.json} lists them; the verifier checks
every coefficient against the Kraus definition above. The factor $96$ is
$\det(D_4)^2$.

For normalized inputs, if $R(\rho)=I_5/5+X_1+X_2+X_3+X_4$, the same channel is
\[
 R\Phi_2(\rho)=\frac{I_5}{5}+\frac{23}{35}X_1+\frac1{10}X_2
                 -\frac{29}{70}X_3-\frac12X_4.
\]
The certificate concerns the homogeneous degree-five polynomial $\Delta$;
no trace normalization or restriction on $X_1,X_3$ is used in the proof.

\subsubsection{The weighted sum-of-squares identity}
Let $S_{\rm a},S_{\rm s},S_{\rm m}$ be the real rational row bases specified in
\path{source_bases.json}. They span the exterior cube, symmetric cube, and
the first-two-factor-symmetric subspace orthogonal to the fully symmetric
cube, respectively. Their dimensions are $20,56,70$.
Define the positive weights
\[
 W_1(H)=H,\qquad
 W_\lambda(H)=S_\lambda H^{\otimes3}S_\lambda^{\mathsf T}
 \quad(\lambda={\rm a,s,m}).
\]
All four weights are positive semidefinite when $H\succeq0$.

For output spin $J$, put
$M_J=\diag\bigl(\binom{2J}0,\ldots,\binom{2J}{2J}\bigr)$.
The rational polynomial matrices $Q_{\lambda,J,a}(H)$ in the final data have
degree two for $\lambda=1$ and degree one for the cubic weights. The identity is
\begin{equation}\label{sos:eq:sos}
 \boxed{\displaystyle
 \Delta(H)=\sum_{\lambda,J}\sum_{a,b}
 G^{(\lambda,J)}_{ab}\,
 \Tr\!\left[W_\lambda(H)
 Q_{\lambda,J,a}(H)^\dagger M_JQ_{\lambda,J,b}(H)\right].}
\end{equation}
The double sums are real because the Gram matrices are real symmetric.
There are 35 Gram blocks, with the following dimensions:
\begin{center}
\begin{tabular}{c|rrrrrrrrr|r}
\toprule
$J$ & $\frac12$&$\frac32$&$\frac52$&$\frac72$&$\frac92$&$\frac{11}2$&$\frac{13}2$&$\frac{15}2$&$\frac{17}2$& Total\\
\midrule
$W_1$ &17&30&36&38&35&26&16&7&2&207\\
$W_{\rm a}$ &5&9&11&11&9&7&4&2&1&59\\
$W_{\rm s}$ &7&12&15&14&11&7&3&1&0&70\\
$W_{\rm m}$ &12&22&27&28&24&17&10&5&1&146\\
\bottomrule
\end{tabular}
\end{center}
Their 5523 independent entries are rational. A real Gram factorization of
the positive Gram blocks turns \eqref{sos:eq:sos} into 482 weighted matrix squares.

Each final polynomial matrix is encoded in
\path{exact_reduced_basis_k.json}: every coefficient is the listed integer
divided by the factor's listed denominator. The matrices $G$ are encoded in
\path{exact_sos_certificate.json} as
\[
 G_k=L_kL_k^{\mathsf T}+E_k.
\]
The entries of $L_k$ are listed integers divided by $10^{12}$; the sparse
corrections $E_k$ are listed exact rational numbers. Upper-triangular entries
are indexed consecutively, block by block and row by row.

\subsubsection{Why the finite identity check is conclusive}\label{sec:sos-identity}
We spell out the polynomial identity test, including the dimension
calculation. No invariant-space dimension is an additional hypothesis.
All coefficients below are rational; the variables $h_{ab}$,
$0\le a,b\le5$, are initially independent. On Hermitian matrices the
formal involution $h_{ab}^*=h_{ba}$ is ordinary complex conjugation.

\paragraph{Exact meaning of the polynomial and Gram data.}
Block $k$ is specified by \path{exact_reduced_basis_k.json}.
Its fields \texttt{kind}, \texttt{J2}, and \texttt{source\_dimension}
choose $\lambda$, $2J$, and the number of columns of $Q$.
For its $a$th polynomial matrix, divide all coefficients by
\texttt{denominators[a]}. A stored entry $[\boldsymbol\mu,c,v]$
in row $r$ contributes
\[
 (Q_{\lambda,J,a})_{r,c}\mathrel{+}=
 \frac{v}{\texttt{denominators[a]}}
 \prod_{h\in\boldsymbol\mu}h_{\lfloor h/6\rfloor,\,h\bmod6}.
\]
Unlisted entries are zero, and repeated entries add. The lists have
length two for $W_1$ and length one for the cubic weights. Thus the
weighted expression has degree five in every block.
The rows of $S_\lambda$ are the integer rows in
\path{source_bases.json}; no orthonormality of those rows is required
for positivity. The field \texttt{rational\_whitening} records
construction history and is not applied a second time to $Q$.
The Gram matrices are
$G_k=L_kL_k^{\mathsf T}+E_k$, where
$L_k=\texttt{L\_integer}/10^{12}$ in
\path{exact_sos_certificate.json}. Its sparse \texttt{corrections}
are the entries of $E_k$, indexed by the upper triangles, block by
block, each triangle in row order; a listed off-diagonal value is
placed symmetrically. These rules specify every entry of the identity.

\paragraph{Positivity without a numerical eigenvalue test.}
For each block the exact rational checks give
\[
 \sum_{i,j}|(L_k^{-1})_{ij}|^2\le2^{16},\qquad
 \max_i\sum_j|(E_k)_{ij}|\le2^{-38}.
\]
Indeed the first bound implies
$\|L_k^{-1}\|_{\rm op}^2\le2^{16}$, so
$L_kL_k^{\mathsf T}\succeq2^{-16}I$; the second bounds the operator
norm of a real symmetric $E_k$. Hence
$G_k\succeq(2^{-16}-2^{-38})I\succ2^{-17}I$.
Factoring $G_k=C_k^{\mathsf T}C_k$ and putting
$Z_s=\sum_a(C_k)_{sa}Q_a$ rewrites its contribution as
\[
 \sum_s\operatorname{Tr}(W_\lambda Z_s^\dagger M_JZ_s)
 =\sum_s\|M_J^{1/2}Z_sW_\lambda^{1/2}\|_{\rm HS}^2\ge0.
\]
This argument uses only positivity of $H$, and covers singular $H$.

\paragraph{Covariance is checked as coefficient identities.}
Put $\widetilde Q_a=Q_aS_\lambda$, with $S_1=I_6$ for the linear
weight. Its source columns are indexed by $c$ for that weight and by
the ordered triples $(c_1,c_2,c_3)$ for the cubic weights. Represent
the latter by $z^{(1)}_{c_1}z^{(2)}_{c_2}z^{(3)}_{c_3}$, where the
three sets of six variables are independent. The stored column index
is $36c_1+6c_2+c_3$. These are coordinates of the full tensor cube:
permutations of the slots are not identified. Identifying the three
sets of variables would erase the exterior and mixed-symmetry source
spaces and would not prove their covariance.
The two derivations on these polynomials are determined by
\begin{align*}
 \mathcal D_{\downarrow}h_{ab}&=(5-a)h_{a+1,b}-b h_{a,b-1},&
 \mathcal D_{\downarrow}z_a^{(s)}&=(5-a)z_{a+1}^{(s)},\\
 \mathcal D_{\uparrow}h_{ab}&=a h_{a-1,b}-(5-b)h_{a,b+1},&
 \mathcal D_{\uparrow}z_a^{(s)}&=a z_{a-1}^{(s)},
\end{align*}
with out-of-range variables zero and the product rule on monomials.
If $q_r$ is lifted row $r$, the checks are
\[
 \operatorname{weight}(q_r)=2J-2r,\qquad
 \mathcal D_{\downarrow}q_r=(2J-r)q_{r+1},\qquad
 \mathcal D_{\uparrow}q_r=rq_{r-1}.
\]
The weight of $h_{ab}$ is $2(b-a)$ and that of $z_c^{(s)}$ is $5-2c$.
These equations are checked coefficient by coefficient for all 482
matrices. They say exactly that each polynomial row transforms as the
indicated spin-$J$ representation. To show explicitly that the weighted
traces are invariant, let $U$ be a physical spin-$5/2$ rotation, set
$B_U=D_5^{-1}UD_5$ and $H_U=B_UHB_U^\dagger$, and write $n=1$ or $3$
for the number of source slots. Integrating the generator identities
gives
\[
 \widetilde Q_a(H_U)B_U^{\otimes n}=B_J(U)\widetilde Q_a(H),
 \qquad B_J(U)^\dagger M_JB_J(U)=M_J.
\]
The second equality follows directly from the ladder coefficients and
the binomial entries of $M_J$. Rotations are products of exponentials
of the generators, so the infinitesimal identities imply these group
identities. By the definition of $W_\lambda$ and cyclicity of trace,
each weighted feature is
\[
 \operatorname{Tr}\!\left[H^{\otimes n}
 \widetilde Q_a(H)^\dagger M_J\widetilde Q_b(H)\right].
\]
Substituting $H_U$ cancels $B_U^{\otimes n}$ and its adjoint against
their inverses in the two transformed $\widetilde Q$ factors; the
displayed metric identity then gives its original value. This proves
invariance on all Hermitian inputs, including mixed inputs in the
exterior and mixed-symmetry blocks.
The determinant side is invariant by the covariance of its two maps.
Both sides are unchanged by $H\mapsto H^{\mathsf T}$: the channel
coefficients are real and respect transpose, and transpose interchanges
the two factors of each Gram feature, leaving the symmetric double sum
unchanged. Thus their difference belongs to the same even invariant space.

\paragraph{Derivation of the invariant-space dimension.}
The complexified matrix space decomposes into one copy of spins
$0,1,\ldots,5$, as proved in Section~\ref{sec:prelim}.
A degree-five polynomial is in its fifth symmetric power. A variable
of rank $\ell$ and weight $m$ contributes degree one, weight $m$,
and parity $\ell\bmod2$. On a rotation-invariant scalar,
transpose has the same action as spin reversal, so this parity is
exactly transpose parity. Count monomials with degree at most five
using the following integer recurrence. Start with $C(0,0,0)=1$
and zero elsewhere. For each pair $(\ell,m)$, $0\le\ell\le5$,
$-\ell\le m\le\ell$, replace $C$ by
\[
 C_{\rm new}(d,w,p)=\sum_{k=0}^d
 C_{\rm old}(d-k,w-km,p-k\ell\bmod2),\qquad 0\le d\le5.
\]
The 36 successive steps give
$C(5,0,0)=22894$ and $C(5,1,0)=22462$.
Every integer-spin irreducible representation contributes one vector
at weight zero and, unless its spin is zero, one at weight one.
Subtracting therefore counts precisely the trivial representations:
\[
 \dim\mathcal I_{5,+}=22894-22462=432.
\]
This is a finite integer computation with the complete recurrence supplied.

\paragraph{Explicit invariant contractions and injectivity.}
There is also an explicit basis test, rather than a numerical sampling
claim. Define spin-$\ell$ moment coordinates by
\[
 t_{\ell,0}(H)=\sum_{a=0}^{5-\ell}\binom{5-\ell}{a}h_{a,a+\ell},
 \qquad t_{\ell,r+1}=
 \frac{\mathcal D_{\downarrow}t_{\ell,r}}{2\ell-r}.
\]
These $36$ coordinates are the tensor-rank decomposition, so the
change of coordinates is invertible. To couple spin-$a$ coordinates
$A$ and spin-$b$ coordinates $B$ to spin $L$, put $k=a+b-L$ and
\[
 (A\mathbin{\otimes_L}B)_0=
 \sum_{i=0}^k(-1)^i\binom ki A_iB_{k-i},\qquad
 (A\mathbin{\otimes_L}B)_{r+1}=
 \frac{\mathcal D_{\downarrow}(A\mathbin{\otimes_L}B)_r}{2L-r}.
\]
The generator relations prove that these are spin-$L$ coordinates.
A stored tuple $(a,b,u,c,d,v,e)$ specifies the degree-five scalar
\[
 I(H)=\sum_{r=0}^{2u}(-1)^r\binom{2u}{r}
 (t_a\mathbin{\otimes_u}t_b)_r
 \bigl((t_c\mathbin{\otimes_v}t_d)\mathbin{\otimes_u}t_e\bigr)_{2u-r}.
\]
The triangular inequalities and even $a+b+c+d+e$ are checked for all
432 stored tuples, so these are elements of $\mathcal I_{5,+}$.
The same file, \path{invariant_grid_certificate.json}, gives 432
Hermitian Gaussian-integer matrices $H_i$ through their real and
imaginary parts. Form the square evaluation matrix $V_{ij}=I_j(H_i)$.
With $p=1000033$ and $i$ represented by $350504$ modulo $p$, exact
elimination gives
\[
 350504^2\equiv-1\pmod p,\qquad \det V\equiv486549\not\equiv0\pmod p.
\]
Primality of $p$ can be checked by trial division up to $1000$; all
rational denominators in this computation are invertible modulo $p$.
A nonzero modular determinant implies a nonzero determinant over
$\mathbb Q(i)$, and hence an invertible evaluation matrix over $\mathbb C$.
The 432 contractions are therefore linearly independent, and the
proved dimension makes them a basis of $\mathcal I_{5,+}$.
Evaluation at the stored points is consequently injective on that space.

Finally substitute every $H_i$ into the two sides of the SOS identity,
using the stated rational encodings and ordinary determinant expansion.
All 432 rational differences are zero. This equality, the covariance
identities, the parity check, and injectivity prove the polynomial
identity on every Hermitian $H$. In particular the test is not an
assumption that finitely many sample points suffice: its injectivity
has just been proved. The embedded verifier implements these exact
steps and its earlier audit is included with the data. No Lean build
is used in this written argument.

The identity and the nonnegative-square argument prove
$\Delta(H)\ge0$ for every $H\succeq0$. Since
$\Delta(H)=\det R_{5,4}\Phi_2^{(5)}(\rho)-\det R_{5,4}(\rho)$,
this completes the fourteenth base case. The hull argument at the
start of this chapter now proves the determinant comparison for every
$N\ge5$, every covariant channel, and every density matrix.

\clearpage\section{The common shifted determinant inequality}\label{sec:shift}
\begin{lemma}\label{lem:shift}
Equation~\eqref{eq:det}, applied to every global input density matrix, implies
\[
 \det(s\id_d+\rho'_t)\ge\det(s\id_d+\rho_t)
 \qquad\text{for all }s\ge0.
\]
\end{lemma}
\begin{proof}
Covariance and irreducibility give $\Phi(\rho_0)=\rho_0$, and
Section~\ref{sec:prelim} gives $R_{N,t}\rho_0=\id_d/d$.
For the full input $\rho_s=(\rho+ds\rho_0)/(1+ds)$, linearity
therefore gives
\[
 R_{N,t}(\rho_s)=\frac{\rho_t+s\id_d}{1+ds},\qquad
 R_{N,t}\bigl(\Phi(\rho_s)\bigr)=\frac{\rho'_t+s\id_d}{1+ds}.
\]
Apply~\eqref{eq:det} to $\rho_s$ and cancel the positive common
factor $(1+ds)^{-d}$.
\end{proof}

\section{From shifted determinants to fidelity and entropy}\label{sec:spectral}
Let $p_i$ and $q_i$ be the eigenvalues of $\sigma=\rho_t$ and $\tau=\rho'_t$.
For $r>0$, Lemma~\ref{lem:shift} gives
\[
 0\le\log\frac{\det(r^2\id_d+\tau)}{\det(r^2\id_d+\sigma)}
 =\sum_{i=1}^d\bigl[\log(1+q_i/r^2)-\log(1+p_i/r^2)\bigr].
\]
To integrate this comparison, fix $x>0$. Integration by parts gives
\[
 \int_0^\infty\log(1+x/r^2)\,dr
  =2x\int_0^\infty\frac{dr}{r^2+x}=\pi\sqrt{x}.
\]
The boundary term $r\log(1+x/r^2)$ tends to zero at infinity by
$\log(1+x/r^2)\le x/r^2$. At zero, write it as
$r\log(r^2+x)-2r\log r$. For $x=0$ the integrand vanishes identically.
Thus each term is integrable, including for singular marginals. Integrating
the finite sum yields $\pi(\Tr\sqrt\tau-\Tr\sqrt\sigma)\ge0$.
Squaring the nonnegative traces and applying~\eqref{eq:F} proves fidelity
monotonicity for every input at each of the four orders.

\subsection{Von Neumann monotonicity from fidelity monotonicity on all inputs}\label{sec:fidelity-implies-entropy}
\begin{lemma}\label{lem:fidelity-implies-entropy}
Fix $1\le t\le N$ and a completely positive, trace-preserving
$\mathrm{SU}(2)$-covariant channel $\Phi$ on spin $N/2$.
If $\A_t^{T,F}(\Phi(\eta))\ge\A_t^{T,F}(\eta)$ for every density matrix
$\eta$ on this spin space, then
$\A_t^{T,\mathrm{vN}}(\Phi(\rho))\ge\A_t^{T,\mathrm{vN}}(\rho)$ for every
density matrix $\rho$ on the same space.
\end{lemma}
\begin{proof}
Fix $\rho$ and put $\sigma=\rho_t$, $\tau=\rho'_t$. For each $s\ge0$
apply the fidelity hypothesis to the full mixture $\rho_s$ in
Lemma~\ref{lem:shift}. Comparing the nonnegative square-root traces in~\eqref{eq:F}
and cancelling the common factor $(1+ds)^{-1/2}$ gives
\[
 \Delta(s)=\Tr\sqrt{\tau+s\id_d}-\Tr\sqrt{\sigma+s\id_d}\ge0 .
\]
The entropy difference is
\begin{equation}\label{eq:entropy-from-fidelity}
 S(\tau)-S(\sigma)=4\lim_{b\to\infty}\int_0^b\Delta(u^2)\,du .
\end{equation}
For an eigenvalue $x\ge0$ and $b>0$, direct integration gives
\[
 \int_0^b\sqrt{u^2+x}\,du
 =\tfrac12\!\left[b\sqrt{b^2+x}+x\log\!\left(b+\sqrt{b^2+x}\right)\right]
   -\tfrac14 x\log x .
\]
For $x=0$ the formula gives $b^2/2$ with $0\log0=0$. As $b\to\infty$,
the bracketed term times $1/2$ equals
$b^2/2+(x/2)\log(2b)+x/4+o(1)$. Indeed,
$b\sqrt{b^2+x}-b^2=x/(1+\sqrt{1+x/b^2})\to x/2$, and
$(b+\sqrt{b^2+x})/(2b)\to1$.
Sum over the eigenvalues of $\tau$ and subtract the sum for $\sigma$.
Equal dimensions cancel the $b^2/2$ terms, and equal traces cancel the
terms linear in $x$, proving~\eqref{eq:entropy-from-fidelity}.
Every integral on its right is nonnegative. Dividing by $\log d>0$
gives the von Neumann inequality in the normalization~\eqref{eq:VN}.
\end{proof}

At $t=1,2,3,4$ the fidelity comparison proved above supplies the hypothesis
of Lemma~\ref{lem:fidelity-implies-entropy}. This proves the second
inequality of Theorem~\ref{thm:main} for the same channel and every input.

\section{Embedded exact data: definitions and checking rules}\label{sec:data}
\subsection{What is contained in this PDF}
The attached archive \texttt{exact-data.zip} is an integral machine-readable
appendix. It contains every rational table and polynomial coefficient used
in the finite checks above, together with the exact checking programs.
Open the PDF's attachments panel to save the archive, then extract it.
All data paths below are relative to that extracted directory. A PDF
viewer that does not expose attachments can still display every written
page; another PDF reader is needed to extract the data. No external
repository, network service, or Lean result is needed to obtain them.

The archive has 113 entries: 110 original audited inputs and programs,
a manifest, an audit record, and a reading guide. Its SHA-256 digest is
\begin{center}\small\ttfamily
79bc6ad1a25aeebf2215beb20c285fc39205\\
21b505ef1bebb48fd1c0ba239801
\end{center}
(The two lines are concatenated to obtain the 64-digit digest.)
\path{MANIFEST.json} gives a byte count and SHA-256 digest for every
original file. Hashes identify the supplied data; the inequalities and
identities, rather than the hashes, are the mathematical certificate.

\begin{center}
\small
\begin{tabular}{@{}p{.21\linewidth}p{.73\linewidth}@{}}
\toprule
Proof step&Exact data and check location\\
\midrule
Order-three hull&\path{determinant_t3/bounds/finite_hull.json},
\path{finite_hull_certificates.json} in the same folder;
\path{determinant_t3/search/audit_reduction.py}.\\[2mm]
Order-three scalar intervals&\path{determinant_t3/bounds/central_region_certificates.json}
and \path{k42_outer_certificates.json} in that folder; the corresponding
\path{verify_central_region.py} and \path{certify_k42_outer.py}.\\[2mm]
Order-three floors&\path{determinant_t3/floors/positive_floors.json}
and \path{verify_positive_floors.py} in that folder.\\[2mm]
Order-four hull&\path{determinant_t4/bounds/exact_hull.json},
\path{continuous_hull_N12.json}, \path{small_case_certificates.json};
\path{determinant_t4/structural/audit_finite_reduction.py}.\\[2mm]
Thirteen floor cases&All thirteen
\path{determinant_t4/bounds/*norm_floor_certificate_N*_K*.json};
\path{verify_norm_floor_certificate.py} in that folder,
\path{determinant_t4/search/audit_norm_floor.py}, and
\path{determinant_t4/structural/audit_k3_compact.py}.\\[2mm]
Exceptional case&\path{determinant_t4/coherent_sos/channel_rational.json};
\path{determinant_t4/exact_cube_sos/} contains
\path{source_bases.json}, \path{exact_reduced_basis_0.json} through
\path{exact_reduced_basis_34.json}, \path{exact_sos_certificate.json},
\path{invariant_grid_certificate.json}, and \path{verify_exact_sos.py}.\\
\bottomrule
\end{tabular}
\end{center}
The remaining included helpers implement rational arithmetic and the
finite constructions defined in the text. The additional exact-kernel
check describes the construction of the SOS candidate; its success
is not needed once the final identity and positivity have been checked.

\subsection{Rational interval arithmetic, with endpoint directions}
A rational $x>0$ has the exact logarithm enclosure obtained as follows.
If $x<1$, use $\log x=-\log(1/x)$ and interchange the signs of the
upper and lower endpoints. Otherwise write $x=2^kR$, $1\le R<2$,
set $z=(R-1)/(R+1)$, and use
\[
 A_M(z)=2\sum_{r=0}^{M-1}\frac{z^{2r+1}}{2r+1},\qquad
 0\le\log R-A_M(z)\le
 \frac{2z^{2M+1}}{(2M+1)(1-z^2)}.
\]
The inequality follows by integrating the geometric series for
$2/(1-z^2)$, or by comparing each omitted denominator with $2M+1$.
Use the same formula at $z=1/3$ for $\log2$ and add $k$ times its
bounds. The original scalar certificates use $M=12$ after range
reduction; the independent order-four audit uses $M=20$.
For $q\ge0$ and a positive integer $B$, put
\[
 h=\frac{\operatorname{isqrt}(\lfloor B^2q\rfloor)}B.
 \quad\text{Then}\quad h^2\le q<(h+1/B)^2.
\]
Here $\operatorname{isqrt}(n)$ is the largest integer with square at
most $n$, computed by integer bisection. The original grid has
$B=2^{40}$ and the independent audit has $B=2^{48}$.
All lower bounds on an input deficit use the lower logarithm endpoint;
all upper bounds on an output deficit use the upper endpoint. The
square-root direction is chosen using the monotonicity proved in the
corresponding scalar lemma. A strict positive rational gap is therefore
a rigorous inequality on the whole specified box.

\subsection{Coverage and the finite predicates}
A list of successful samples would not prove a statement on a continuum.
The certificates instead specify complete partitions. To check a
partition of a root rectangle, accept a leaf only if its endpoints equal
the current rectangle; otherwise require a split in a single coordinate,
strictly between its two endpoints, with one child on each side and
unchanged other coordinates. Apply the same rule recursively to both
children, requiring every stored leaf to be used exactly once.
Induction on this binary tree proves that the union is exactly the root
rectangle and interiors do not overlap. Shared boundaries are included
because the predicates use closed intervals. The scalar-box checkers
reconstruct this tree from the stored rational endpoints.

A scalar leaf is accepted only after one of the following checks:
\begin{enumerate}
\item A necessary physical inequality fails throughout the leaf. Each
positive image used for that inequality has first been constructed by
the explicit Kraus or positive-map formulas and its multiplier identity
checked exactly.
\item The input is forced to have determinant zero, in which case output
positivity is enough.
\item A positive spectral floor has been derived from an exact rational
nonnegative combination of those physical images, with the remaining
multipole terms bounded by the proved operator norms. Its output-deficit
upper bound is no larger than a proved input-deficit lower bound over
the whole leaf. Linear inequalities use endpoint minima; ratios are
checked at the vertices of their linear-fractional domains, including
new vertices cut out by the stated comparison planes.
\end{enumerate}
An invalid denominator, unavailable auxiliary map, nonpositive required
floor, or unresolved comparison causes rejection; a stored success flag
alone is never a reason for acceptance.

For an all-spin hull leaf, substitute its affine box coordinates into
the facet polynomial and expand exactly. Convert its power coefficients
to Bernstein coefficients using the formula in the hull chapters.
The identity
\[
 x^i=\sum_{k=i}^n\frac{\binom ki}{\binom ni}
 \binom nk x^k(1-x)^{n-k}
\]
follows by factoring $x^i$ and expanding $(x+(1-x))^{n-i}$.
Taking its product in the two coordinates proves that conversion
formula. Every Bernstein basis function is nonnegative on the unit
square; their sum is one. Nonnegative rational coefficients therefore
prove the whole leaf, including its boundary. The complete partition
then proves the facet on the entire continuous domain. Exact finite
hull enumeration and the small-spin convex weights establish that
these facets are sufficient and that every extension used is available.

\subsection{Scope of the previously completed check}
The attached \path{prior-exact-audit.json} records a successful independent
exact-Python run completed on 16 November 2026. During this PDF edit,
all 110 original input and program files were compared byte for byte
through their hashes with that audited manifest; every comparison matched.
The scalar proof printed here uses 55\,003 ordinary order-four leaves
and 1\,184 compact endpoint-region leaves, a total of 56\,187. The exceptional proof here uses the proved 432-dimensional invariant
grid, whose dimension and injectivity were derived above.

The full checking entry point is \path{verify_all.py}; its dependencies
are listed in \path{requirements.txt}. Reproduction requires exact
integer/rational arithmetic, with the Clebsch--Gordan constructions
implemented symbolically. Floating-point search, optimization, and
numerical eigenvalue estimates are not accepted proof steps. The
checking programs are supplied to make all the finite arithmetic
reviewable; the mathematical reasons why their checks suffice are
written out in this document.

\subsection{Written-proof audit of 16 November 2026}
A fresh run extracted the exact archive from the PDF, matched all 110
original file hashes, and passed the order-three suite, the order-four
spectral-floor suite, the exceptional weighted SOS and positivity check,
and the kernel-constraint check. Additional exact checks confirmed the
36 independent moment coordinates, invariant-space dimension $432$, and
all 24 symbolic band identities for arbitrary $N$.

The covariance argument retains three independent ordered
source slots and explicitly proves weighted-trace invariance. This is
essential for the exterior and mixed-symmetry blocks. The certificate
already retained these slots, so its exact data are unchanged.
These exact-arithmetic checks are independent of the Lean certificates below.

\clearpage\section{Lean certificates}\label{sec:lean}
The public theorems use $N=k+t$ with $k\ge0$. As in
Theorem~\ref{thm:main}, the channel is completely positive,
trace preserving and covariant under the spin representation of
$\mathrm{SU}(2)$, and the full input is positive semidefinite of trace one. The reduced state is the
normalized Dicke marginal. The two totals are those of~\eqref{eq:F}
and~\eqref{eq:VN}.

\paragraph{Common argument.}
For a fixed channel, \path{LowOrder.physical_fidelity_of_determinant}
turns the determinant inequality on every full input into fidelity
monotonicity on every full input. It first applies
\path{shifted_determinant_of_density_determinant} to isotropic mixtures
(Section~\ref{sec:shift}), then takes the square-root trace comparison
\path{matrix_sqrt_trace_mono_of_shifted_det}
(Section~\ref{sec:spectral}). The theorem
\path{LowOrder.physical_vonNeumann_of_fidelity} takes this all-input
fidelity result as its hypothesis and yields entropy monotonicity.
The combined theorem \path{LowOrder.physical_totals_of_determinant}
in \path{PhysicalUniformBridge.lean} composes the two steps.
The \href{https://github.com/aenictusGitHub/mixed-spin-anticoherence/tree/main/fidelity-implies-entropy}{general proof and Lean certificate}
give the second step for arbitrary $1\le t\le N$.

\paragraph{Order-specific arguments.}
\begin{itemize}
\item At $t=1$, \path{lean/UniformT1.lean} uses positivity and unitality
on $jI\pm J_z$ to bound the dipole multiplier by $|f_1|\le1$.
The determinant difference is $(1-f_1^2)(1/4-\det\rho_1)$.
At $t=2$, \path{lean/ChannelSimplexT2Uniform.lean} proves the determinant
comparison in the physical marginal coordinates. In
\path{lean/UniformFidelity.lean}, \path{physical_fidelity_t1} and
\path{physical_fidelity_t2} supply the hypotheses of
\path{physical_vonNeumann_t1} and \path{physical_vonNeumann_t2}, respectively.
\item At $t=3$, \path{lean-t3/T3PhysicalDeterminant.lean}
proves \path{physical_determinant_t3} using the four finite cases,
their physical constraints and the available-extension hull.
The two conclusions are \path{physical_fidelity_t3} and
\path{physical_vonNeumann_t3} in \path{lean-t3/UniformT3Certificate.lean}.
\item At $t=4$, \path{lean-t4/T4PhysicalDeterminant.lean}
proves \path{physical_determinant_t4} from the fourteen base cases
and the all-spin hull. For the exceptional $(N,K)=(5,2)$ case,
\path{lean-t4/T4ExceptionalFinal.lean} uses 35 exact Gram certificates
and an invariant-polynomial argument with 687 seed coefficients.
The two conclusions are \path{physical_fidelity_t4} and
\path{physical_vonNeumann_t4} in \path{lean-t4/UniformT4Certificate.lean}.
\end{itemize}

\paragraph{Verification.}
Each package pins Lean 4.24.0 and its Mathlib revision. The
\href{https://github.com/aenictusGitHub/mixed-spin-anticoherence/tree/main/su2-low-order-monotonicity/lean}{order-one and order-two},
\href{https://github.com/aenictusGitHub/mixed-spin-anticoherence/tree/main/su2-low-order-monotonicity/lean-t3}{order-three}, and
\href{https://github.com/aenictusGitHub/mixed-spin-anticoherence/tree/main/su2-low-order-monotonicity/lean-t4}{order-four}
packages give the reproduction commands, checked sources and axiom reports.

\end{document}
